\documentclass[11pt]{article}
\usepackage[T1]{fontenc}
\usepackage[utf8]{inputenc}
\usepackage{amsmath,amssymb,mathrsfs}
\usepackage[a4paper,margin=28mm]{geometry}
\usepackage[hidelinks]{hyperref}
\newcommand{\dd}{\mathrm{d}}
\newcommand{\id}{\mathbf 1}
\setlength{\parindent}{0pt}
\setlength{\parskip}{0.6em}
\emergencystretch=1.5em
\title{Fourth-Order Tests for the Inherited C8 Wilson Insertion}
\author{}
\date{Working attempt, 22 October 2026}
\begin{document}
\maketitle

\section{A Fourth-Order Test on the Inherited C8 Insertion}
\label{sec:c8-fourth-test}

This section records the calculation authorised on 22 October 2026:
``Yes, attempt the fourth-order extension for the actual C8 Wilson
insertion, retaining the full inherited law and physical denominator.
Don't use Codex subagents.'' The target is a quantitative estimate for
the source-containing kernels WC-C8.6, with first use in the boundary
bracket LFR2 (1.103), the connected term WC-C7.5, and ultimately WC.1.
The finite first-step signed defect in V3, equation
\texttt{eq:c4-1-1}, is a related caller; it is not a replacement for
the inherited C8 target. The calculation below tests the positive-square
step of the fourth-order NLSM argument~\cite{OpenAIRotation}. It does not assume that the
NLSM noncommuting-rotation identity already holds for these kernels.

Whenever two loops are considered, the target loops are disjoint.
A component meeting both loops means a connected integration
component or interaction cluster reaching both disjoint source
supports; it does not mean an intersection of the contours.
Such connecting components are retained under their common
Yang--Mills law. Cancellation of components not connecting the
sources is a statement about the connected covariance, not
about the uncentred joint expectation.

\subsection{The Literal Operand and a Defined Four-Source Test}

The authoritative definitions are WC-C8.1--WC-C8.7 in
\texttt{WILSON\_RG\_COMPARISON.txt}. Fix a finite cutoff and complete
outer data $\eta$ in LFR2 (1.103), on the stated domain where the
component normalisers and the vacuum bracket denominator are nonzero.
The character is $\chi=\operatorname{Tr}R/d$ for a nontrivial irreducible
unitary representation $R$. The scalar $Z\neq 0$ is the already prescribed
source renormalisation. For a nonempty marked set with $m$ entries,
the bare scalar kernel is
\begin{equation}\label{eq:c8-fourth-bare}
    b_J(H)=\frac{1}{Zd}\operatorname{Tr}
        \prod_{r=1}^{m}\left[(R(Q_{a_r})-\id)R(H_r)\right]\;.
\end{equation}
Here $H\in G^m$ are independent auxiliary closed holonomies,
$Q_a=T_aF_aV_a^{-1}T_a^{-1}$ are the closed relative loops of WC-C8.2,
and all products retain their cyclic order. These are scalar closed
contractions. Every factor is conjugated at the same base under the
induced gauge action, so the trace is invariant before any estimate.

Let $b=\bar b_{J,\sigma}(\eta;H)$ be an actual inherited labelled
marker, including its earlier numerator and touching-denominator labels.
Its unlabelled sum is the chronological quotient WC-C8.3a, with the
replacement rule WC-C8.3b where applicable. Earlier integrated variables
are not reconstructed from $\eta$. Put
\begin{equation}\label{eq:c8-fourth-bracket}
    L_\eta=\prod_i\frac{T'_k(Y_i)}{T'_k(Y_i)1},\qquad
    B_\eta=\sum_a\beta_a,\qquad
    D_\eta=L_\eta[e^{B_\eta}],\qquad
    \Lambda_\eta[F]=\frac{L_\eta[Fe^{B_\eta}]}{D_\eta}\;.
\end{equation}
The operations and their variables are precisely LFR2 (1.71),
pp.~378--379, acting on the last exponential of (1.72), and then
(1.103), p.~391. Each $\beta_a$ is an input boundary term. In
particular $b$, $B_\eta$ and the tests below remain inside all component
operations that read them. A normalised functional on the complex
extension is not declared positive. LFR2 (1.73)--(1.75) supplies a
positive real reference operation and a finite-variation comparison,
not positivity of the complex extension or of selected activities.

Choose finitely many actual bounded input boundary terms and let
$p_i=\operatorname{Re}\beta_{\alpha_i}$ and
$q_i=\operatorname{Re}\beta_{\gamma_i}$ on the real integration
variables; $i$ ranges over a fixed finite set. These are real,
gauge-invariant scalar functions independent of the auxiliary $H$.
This is a diagnostic real source deformation; no holomorphic extension
of a real-part operation in the background is asserted. For a pair
$(i,j)$ set $x_1=p_i$, $x_2=q_i$, $x_3=p_j$, $x_4=q_j$ and define
\begin{equation}\label{eq:c8-fourth-defined-test}
    \begin{aligned}
        N^{ij}(s;H)&=L_\eta\left[b(H)e^{B_\eta+\sum_{r=1}^{4}s_rx_r}\right],\\
        D^{ij}(s)&=L_\eta\left[e^{B_\eta+\sum_{r=1}^{4}s_rx_r}\right],\qquad
        F^{ij}(s;H)=\frac{N^{ij}(s;H)}{D^{ij}(s)}\;.
    \end{aligned}
\end{equation}
The four real parameters are independent, also when some tests coincide.
At zero this is the original C8 quotient. The deformation multiplies
the actual boundary integrand before the old integrations; it does not
replace those integrations. Boundedness of the tests and absolute
integrability of the original sourced integrand justify its derivatives
at fixed data by dominated convergence. Nonvanishing of $D_\eta$
gives a neighbourhood of zero for the quotient, with no uniform radius
claimed. No sourced expectation is used as a denominator.

For the full physical version, let $\mathcal O_\eta$ denote exactly
the remaining sums and integrations of LFR2 (1.72)/(1.104), and include
$A_\eta=\prod_i(T'_k(Y_i)1)$ among its multiplicative factors. The
notation suppresses no additional factor: the characteristic functions,
older operations, actions, Haar factors, gauges and vacuum scalars in
those formulas remain in $\mathcal O_\eta$. Define
\begin{equation}\label{eq:c8-fourth-physical}
    N_{\mathrm{phys}}(s)=\mathcal O_\eta[A_\eta N_\eta(s)],\qquad
    D_{\mathrm{phys}}(s)=\mathcal O_\eta[A_\eta D_\eta(s)],\qquad
    D_{\mathrm{phys}}(0)=D_h\;.
\end{equation}
All source and history sums required by the physical numerator are
included when this formula is used. Auxiliary holonomies are evaluated
at their physical values inside these operations. At the first step
$D_h=\mathcal Z_0$, the original full physical partition function. Differentiation
after the outer operations additionally requires integrable derivative
majorants for those operations; the fixed-data calculation alone does
not prove them. The quotient rule below applies to the physical pair
whenever those derivatives exist. It never cancels $D_h$ against a
conditional reference normaliser.

\subsection{Four Derivatives and the Denominator Terms}

For a subset $S$ of four labelled directions write $N_S$, $D_S$, and
$F_S$ for their mixed derivatives at zero. Differentiating $N=DF$
gives, for every nonempty $S$,
\begin{equation}\label{eq:c8-fourth-quotient}
    F_S=\frac{N_S-\sum_{\varnothing\neq A\subseteq S}
        D_AF_{S\setminus A}}{D_\varnothing}\;.
\end{equation}
Each labelled subset occurs once. This formula also governs every
earlier chronological quotient WC-C8.3a; it retains derivatives of
older densities and their normalisers when those data vary. Thus a
source derivative is not replaced by a derivative of the bare marker
after its old variables have been eliminated.

For comparison with a nonlinear rotation deformation, write its smooth
relative exponent as $h(s)$, with $h(0)=0$, and write $b(s)$ for the
actual differentiated insertion. This notation is used only when a
common integration domain and dominated derivatives have been justified.
Let $h_A=\partial_Ah(0)$, $b_A=\partial_Ab(0)$, and let $\Pi(S)$ be
the set of partitions of $S$ into nonempty blocks. Set
\begin{equation}\label{eq:c8-fourth-bell}
    P_\varnothing=1,\qquad
    P_S=\sum_{\pi\in\Pi(S)}\prod_{A\in\pi}h_A\;.
\end{equation}
The product rule and differentiation of the exponential give
$D_S/D_\varnothing=\Lambda_\eta[P_S]$ and
$N_S/D_\varnothing=\sum_{A\subseteq S}
\Lambda_\eta[b_AP_{S\setminus A}]$. In particular,
\begin{equation}\label{eq:c8-fourth-centred}
    F_S=\Lambda_\eta[(b-F_\varnothing)P_S]
        +\sum_{\varnothing\neq A\subseteq S}
            \Lambda_\eta[b_AP_{S\setminus A}]
        -\sum_{\varnothing\neq A\subsetneq S}
            \Lambda_\eta[P_A]F_{S\setminus A}\;.
\end{equation}
For four directions $P_S$ has the fifteen set-partition terms: one
$h_1h_2h_3h_4$, six $h_{ab}h_ch_d$, three $h_{ab}h_{cd}$, four
$h_{abc}h_d$, and one $h_{1234}$. Formula
\eqref{eq:c8-fourth-centred} also includes fifteen nonempty insertion
derivative terms and fourteen proper-subset normaliser terms. Actual
hard-cutoff or constraint variations must first be retained as moving
domain or surface terms; they are not made smooth by writing $h$.
No such variation occurs in \eqref{eq:c8-fourth-defined-test}.

For the defined test, $h(s)=\sum_rs_rx_r$ and $b(s)=b$ exactly.
Set $Y=\sum_ip_iq_i$ and define the retained denominator contribution
\begin{equation}\label{eq:c8-fourth-denominator}
    \mathcal D_b=\sum_{i,j}\sum_{\varnothing\neq A\subsetneq[4]}
        \Lambda_\eta\left[\prod_{r\in A}x_r^{ij}\right]
        F^{ij}_{[4]\setminus A}\;,
    \qquad [4]=\{1,2,3,4\}\;.
\end{equation}
The unnormalised and normalised four-derivative contractions are,
respectively,
\begin{align}
    \sum_{i,j}\frac{N^{ij}_{[4]}}{D_\eta}
        &=\Lambda_\eta[bY^2]\;,
        \label{eq:c8-fourth-square}\\
    \sum_{i,j}F^{ij}_{[4]}+\mathcal D_b
        &=\Lambda_\eta[(b-F_\varnothing)Y^2]\;.
        \label{eq:c8-fourth-connected-square}
\end{align}
These are calculations under the actual inherited functional, not under
a Gaussian law. At a positive real instance, the second right-hand side
is the connected correlation of $b$ and $Y^2$. If its second moments
are finite, Cauchy--Schwarz supplies only
\begin{equation}\label{eq:c8-fourth-cs}
    \left|\sum_{i,j}F^{ij}_{[4]}+\mathcal D_b\right|
        \leq\Lambda_\eta[|b-F_\varnothing|^2]^{1/2}
        \Lambda_\eta[(Y^2-\Lambda_\eta[Y^2])^2]^{1/2}\;.
\end{equation}
There is no small history, coupling or support factor in this inequality.
For complex functionals even this probabilistic form is not asserted.

The positive square in the NLSM proof arises more generally from the
four terms grouped as
$(h_1h_2+h_{12})(h_3h_4+h_{34})$. In the conjugate paired directions
of that proof their sum is $|\mathcal S|^2$. Multiplying the original
integrand by $b$ gives $b|\mathcal S|^2$, and normalising produces
$(b-F_\varnothing)|\mathcal S|^2$, in addition to the insertion
derivatives and proper-subset terms of \eqref{eq:c8-fourth-centred}.
This is the precise term for which the NLSM lower-bound argument would
need a new estimate.

\subsection{An Obstruction in the Actual Holonomy Carrier}

For a nontrivial irreducible $R$, Haar averaging gives
$\int_G R(H)\,\dd H=0$: the average is the projection onto invariant
vectors, and the representation has none. Applied inside the closed
trace in \eqref{eq:c8-fourth-bare}, this proves, in any one auxiliary
slot with the other slots fixed,
\begin{equation}\label{eq:c8-fourth-haar-zero}
    \int_G b_J(H_1,\ldots,H_r,\ldots,H_m)\,\dd H_r=0\;.
\end{equation}
No open factor is separately estimated. Chronological integration,
division by the vacuum normalisers, selection, and denominator
completion WC-C8.6 are linear in the insertion and independent of
$H$. Thus the same equality holds for $\bar b_{J,\sigma}$, for each
completed kernel on its finite-cutoff convergence domain, and for
\begin{equation}\label{eq:c8-fourth-zero-square}
    P(H)=\Lambda_\eta[b(H)|\mathcal S|^2],\qquad
    P_{\mathrm c}(H)=\Lambda_\eta[(b(H)-\Lambda_\eta[b(H)])
        |\mathcal S|^2]\;,
\end{equation}
whenever the integrals exist and $\mathcal S$ is independent of $H$.
The first calculation above has $\mathcal S=Y$. Finite-dimensionality
of the representation makes these functions continuous in $H$.
Haar integration here tests a function on the existing auxiliary
carrier; it is not an extra physical integration or a new probability
law for retained fields.

Consequently, for any fixed phase $\vartheta$, if
$\operatorname{Re}(e^{-i\vartheta}P(H))$ is nonnegative for every
auxiliary $H$, then it is identically zero. Indeed its integral in one
slot is zero, and a continuous nonnegative function with zero integral
against Haar measure is zero. The same conclusion holds for $P_{\mathrm c}$.
If such a real part is nonzero, it takes both signs. This rules out
discarding a nonzero Wilson-weighted square as uniformly nonnegative
on this carrier. It does not assert that every particular family
produces a nonzero $P$, or exclude a bound on its negative part.

The one-mark test is particularly explicit. At $H_1=\id$,
\begin{equation}\label{eq:c8-fourth-one-mark}
    \operatorname{Re}(Zb_{\{a\}}(\id))
        =\operatorname{Re}\chi(Q_a)-1\leq0\;.
\end{equation}
This holds pointwise for the actual relative loop, without changing its
law. For the fundamental representation of $SU(2)$ and
$Q_a=\operatorname{diag}(e^{i\theta},e^{-i\theta})$ it equals
$\cos\theta-1$, strictly negative for $0<|\theta|<2\pi$.
This is a test of the algebraic source carrier, not a claim that an
arbitrary prescribed relative holonomy has been realised by a given
conditioned history. At $Q_a=\id$ the marked expression is zero, as
required by the pure-gauge test. A flat background does not imply
that its integrated fluctuating relative loops are identity.

For the single-germ carrier WC-C8.7, replacing the last argument by
$L_{\mathrm{loc}}H$ preserves zero Haar average by translation
invariance. The prescribed additive scalar $c\chi(H)$ also has zero
Haar average. Subtracting it therefore does not make a nonzero residual
nonnegative. A pair of Wilson insertions retains separate auxiliary
lists, so averaging a slot of either nontrivial character has the same
effect. No independent law for the two loops is used.

\subsection{Caller Replay and Scope of the Failed Extension}

The obstruction concerns a proposed uniform sign step on the C8
carrier. After evaluating $H=H^{\mathrm{phys}}$, the auxiliary argument
depends on the actual retained fields; the zero-Haar-average argument
does not prove a negative physical correlation. Similarly, positivity
of a fully resummed physical law does not imply positivity of an
individual selected kernel. Summing the signed terms before estimating
is still allowed, but requires control of that sum.

Ba{\l}aban's corresponding calculation in CRE (3.44)--(3.47),
pp.~277--278, puts the insertion in $H'(Z_0)$, cancels only the
identical disconnected vacuum factors, and retains the exponential
of clusters touching $Z_0$. Those clusters are precisely the
denominator completion in WC-C8.6. They remain in
\eqref{eq:c8-fourth-quotient}; differentiating a numerator alone would
miss them. LFR2 p.~391 explicitly provides no tree decay inside the
fixed inherited $Y_i$. Neither source supplies a bound for
\eqref{eq:c8-fourth-connected-square} with the required marked gain.

The authorised attempt therefore yields the exact fourth-order
identities \eqref{eq:c8-fourth-square}--\eqref{eq:c8-fourth-connected-square}
and the carrier obstruction \eqref{eq:c8-fourth-haar-zero}.
The unchanged target WC.1 is not proved. No new small activity bound,
ultraviolet uniformity, noncommuting YM rotation identity or all-scale
correlation inequality is asserted. The bounded boundary-source test
directly tests the positivity step of direct transplantation; it is not a full execution
of the noncommuting field-rotation programme. The first unresolved
estimate is the negative/connected part of the source-weighted square,
together with the insertion and denominator terms in
\eqref{eq:c8-fourth-centred}, in the actual marked norm. A subsequent
attempt must bound these terms before claiming a fourth-order gain.

\subsection{Testing Positive Combinations Without a New Denominator}

The continuation authorised on 22 October remains in this separate file.
Consider a genuinely positive real instance of the unexpanded physical
law, normalised by its original $D_h$, or a positive real conditional
operation with its own original vacuum normaliser. Denote its expectation
by $\mathbb E$ in this subsection. No selected complex activity is
declared positive. Let $a$ be a real component of the unlabelled C8
insertion, and suppose $|a|\leq M$. At the primitive level one can take
$a=\operatorname{Re}(e^{-i\vartheta}Zb_J)$ and $M=2^m$: expanding the
$m$ differences gives $2^m$ normalised characters of closed unitary
words, each of modulus at most one. A positive normalised chronological
lift preserves this bound; an individual signed selection need not.

For the nonnegative square $T=|\mathcal S|^2$, the two nonnegative
numerator insertions $M\pm a$ give exactly
\begin{equation}\label{eq:c8-fourth-positive-shifts}
    \mathbb E[(M\pm a)T]=M\mathbb E[T]\pm\mathbb E[aT]\geq 0\;.
\end{equation}
They use the same denominator as the original expectation. Consequently,
with $a_0=\mathbb E[a]$,
\begin{equation}\label{eq:c8-fourth-shift-cost}
    -(M+a_0)\mathbb E[T]\leq
        \operatorname{Cov}(a,T)
        \leq(M-a_0)\mathbb E[T]\;.
\end{equation}
Thus recovering the signed part from the two positive numerators
retains the full vacuum square $\mathbb E[T]$. The NLSM argument
discards that square in a lower bound; it supplies no upper estimate
for its size. This calculation therefore does not recover the desired
fourth-order gain. It does not change the physical normaliser or use
the positive combinations as a new inductive carrier.

\subsection{A Relative-Holonomy Bound for the Connected Square}
\label{sec:c8-fourth-deficit}

We now estimate the actual C8 marker inside the square term. For each
marked relative loop define the nonnegative gauge-invariant scalar
\begin{equation}\label{eq:c8-fourth-deficit}
    q_r=1-\operatorname{Re}\chi(Q_{a_r})\;.
\end{equation}
The next inequality concerns complete closed contractions on the real
compact-group domain, with the same nontrivial irreducible unitary
representation as \eqref{eq:c8-fourth-bare}. For every $H\in G^m$ and
every marked index $r$, it states
\begin{equation}\label{eq:c8-fourth-pointwise-deficit}
    |Zb_J(H)|^2\leq 2^{2m-1}q_r\;.
\end{equation}
To prove it, first let $A=R(Q)$ and let $U$ be any unitary product of
the other based closed holonomies. The Hilbert--Schmidt
Cauchy--Schwarz identity applied within the trace gives
\begin{equation}\label{eq:c8-fourth-closed-cauchy}
    \left|\frac{1}{d}\operatorname{Tr}[(A-\id)U]\right|^2
        \leq\frac{1}{d}\operatorname{Tr}[(A-\id)(A^*-\id)]
             \frac{1}{d}\operatorname{Tr}[U^*U]
        =2(1-\operatorname{Re}\chi(Q))\;.
\end{equation}
Both sides are scalars invariant under simultaneous conjugation at the
base. No open transport is assigned a separate estimate or activity.
In \eqref{eq:c8-fourth-bare}, expand the other $m-1$ differences,
without expanding the chosen $R(Q_{a_r})-\id$. There are $2^{m-1}$
closed traces. Cyclicity places the chosen difference first in each;
the remaining product is unitary. Each trace has magnitude at most
$\sqrt{2q_r}$ by \eqref{eq:c8-fourth-closed-cauchy}. The triangle
inequality for these scalar traces and then squaring prove
\eqref{eq:c8-fourth-pointwise-deficit}. This explicitly pays the
factor $2^{2m-1}$; no cancellation among marked contacts is credited.

For the two-loop C8 insertion, use its product of two complete traces
and set $m=m_1+m_2>0$ and $Z=Z_1Z_2$. Choose the marked index in
either loop. Expanding every other difference gives the same factor
$2^{m-1}$; any remaining unmarked trace has modulus at most one.
Thus \eqref{eq:c8-fourth-pointwise-deficit} also holds for that joint
marker. Its subsequent lift must be the joint chronological lift,
not the product of two separately averaged sources. All covariance
bounds below then use the same joint physical law.

There is an exact inherited version on a positive real chronological
branch. Start $q_{r,0}$ with \eqref{eq:c8-fourth-deficit}, and lift it
through each already performed elimination by the same actual rule
as the unlabelled source:
\begin{equation}\label{eq:c8-fourth-deficit-lift}
    q_{r,\mathrm{next}}(y)=
        \frac{\sum_\lambda\int
            w_\lambda(z,y)q_{r,\mathrm{previous}}(z,y)\,\dd z}
             {I(y)}\;.
\end{equation}
The variables, history labels, base measures, cutoffs and factors are
exactly WC-C8.3a. This formula is used here only when the unexpanded
weights are nonnegative and $I(y)>0$. For an actual replacement
WC-C8.3b with nonnegative $r$ and $\sigma$, its conditional quotient
uses $r/\int r$ for both $b$ and $q_r$; the vacuum output
$\sigma(\int r)/J$ remains in the outer physical weight. No such
positivity statement is made for an individual signed selection or
a complex background extension.

Conditional Cauchy--Schwarz, applied to that same normalised integral,
propagates \eqref{eq:c8-fourth-pointwise-deficit}. Consequently the
unlabelled inherited marker and inherited deficit obey
\begin{equation}\label{eq:c8-fourth-inherited-deficit}
    |Z\bar b_J(\eta;H)|^2
        \leq2^{2m-1}\bar q_r(\eta)\;.
\end{equation}
The proof is induction over the actual finite elimination chronology:
$|\mathbb E_y[Zb]|^2\leq\mathbb E_y[|Zb|^2]
\leq2^{2m-1}\mathbb E_y[q_r]$ at each operation. Postponed variables
stay in $y$. The scalar $\bar q_r$ is only a certificate for this
inequality; it is not substituted for the Wilson insertion or
assigned a new activity norm. All supports and old dependence remain.

Let $\nu$ now be the actual positive real current bracket, including
its boundary exponential, with the unlabelled $\bar b_J$ inside it.
Assume the displayed moments exist, and let $T$ be a real scalar
function of its actual integration variables. Writing the covariance
as $\nu[\bar b_J(T-\nu[T])]$, Cauchy--Schwarz and
\eqref{eq:c8-fourth-inherited-deficit} prove
\begin{equation}\label{eq:c8-fourth-deficit-covariance}
    |\operatorname{Cov}_\nu(\bar b_J(H),T)|
        \leq\frac{2^{m-1/2}}{|Z|}
            \left(\nu[\bar q_r]\,
                \nu[(T-\nu[T])^2]\right)^{1/2}\;.
\end{equation}
This holds uniformly over independent auxiliary $H$. Because
\eqref{eq:c8-fourth-inherited-deficit} is pointwise and uniform in
$H$, it also applies when $H$ is a function of retained variables,
provided its substitution precedes every integration reading it,
as required by C8. The same argument can be applied to the full
unexpanded positive physical law, normalised by $D_h$, if the required
source lift and moments have been justified there. It does not license
a passage through the physical outer integrations without those checks.

Taking $T=Y^2$ in the actual four-boundary-source test yields
\begin{equation}\label{eq:c8-fourth-deficit-replay}
    \left|\sum_{i,j}F^{ij}_{[4]}+\mathcal D_{\bar b_J}\right|
        \leq\frac{2^{m-1/2}}{|Z|}
            \left(\nu[\bar q_r]\,
                \nu[(Y^2-\nu[Y^2])^2]\right)^{1/2}\;.
\end{equation}
The $F^{ij}$ on this line use the unlabelled $\bar b_J$, after summing
its earlier labels, rather than a single $\bar b_{J,\sigma}$. This is
a bound for the connected square together with the fourteen retained
denominator terms, not for $\sum F^{ij}_{[4]}$ alone. No estimate for
those fourteen terms is inferred from it. The second factor contains
eighth-degree moments of the boundary scores, because $Y$ is quadratic
in those scores. Their boundedness does give a finite-data estimate.
Define, using the essential suprema under this actual real law,
\begin{equation}\label{eq:c8-fourth-score-size}
    \mathcal M=\sum_i\|p_i\|_{L^\infty(\nu)}
                          \|q_i\|_{L^\infty(\nu)}\;.
\end{equation}
Then $0\leq Y^2\leq\mathcal M^2$ almost surely. Writing
$t_0=\nu[Y^2]$, the pointwise inequality $Y^4\leq\mathcal M^2Y^2$
implies
\begin{equation}\label{eq:c8-fourth-score-variance}
    \nu[(Y^2-t_0)^2]\leq
        t_0(\mathcal M^2-t_0)\leq\frac{\mathcal M^4}{4}\;.
\end{equation}
Consequently the explicit finite-data bound is
\begin{equation}\label{eq:c8-fourth-finite-bound}
    \left|\sum_{i,j}F^{ij}_{[4]}+\mathcal D_{\bar b_J}\right|
        \leq\frac{2^{m-3/2}}{|Z|}\mathcal M^2
            \nu[\bar q_r]^{1/2}\;.
\end{equation}
No independence or replacement distribution enters this estimate.
Since $0\leq\bar q_r\leq2$ on the positive chronology, the right-hand
side is finite at the fixed data. However, $\mathcal M$ uses an
absolute sum over the family. If there are $n$ tests and all the
individual score bounds are $\delta$, its bound is $n\delta^2$,
so \eqref{eq:c8-fourth-finite-bound} has an $n^2\delta^4$ cost.
This bound supplies no favourable multiscale family cancellation.
Nor does it supply uniform control of $\mathcal M(\eta)$ across the
physical histories, or smallness of the actual inherited deficit.

\subsection{The Square-Root Deficit Loss Is Sharp on This Carrier}

The one-mark case tests whether scalar subtraction could improve
\eqref{eq:c8-fourth-pointwise-deficit} to a bound linear in $q_r$ for
the magnitude of the marker. In the fundamental $SU(2)$ representation,
take the based closed holonomies
\begin{equation}\label{eq:c8-fourth-sharp-holonomies}
    Q_\theta=\operatorname{diag}(e^{i\theta},e^{-i\theta}),
    \qquad H_*=\operatorname{diag}(i,-i)\;.
\end{equation}
Then $\chi(H_*)=0$, $Zb_{\{a\}}(H_*)=-\sin\theta$ and
$q=1-\cos\theta$. For $0<|\theta|<\pi$, the ratio is
\begin{equation}\label{eq:c8-fourth-sharp-ratio}
    \frac{|Zb_{\{a\}}(H_*)|^2}{2q}
        =\frac{1+\cos\theta}{2}\;.
\end{equation}
It tends to one as $\theta$ tends to zero.
In particular $|Zb_{\{a\}}(H_*)|/q$ diverges as
$\theta\longrightarrow0$. Since $\chi(H_*)=0$, subtracting any
scalar multiple of $\chi(H)$ cannot remove this term at $H_*$.
Thus a uniform-in-$H$ estimate $|Zb-c\chi(H)|\leq Cq$ is false on
the algebraic C8 carrier, with a constant independent of $Q$ near
identity. The weaker square-root bound is sharp there. This is not a
claim that every such prescribed $Q_\theta$ is realised by each fixed
RG history, nor a counterexample after physical integration. It tests
the proposed pointwise mechanism on its stated carrier. At
$\theta=0$ the marker vanishes, preserving the pure-gauge test.

The continuation has therefore produced
\eqref{eq:c8-fourth-deficit-covariance}, its exact fourth-order
replay \eqref{eq:c8-fourth-deficit-replay}, and the finite-data bound
\eqref{eq:c8-fourth-finite-bound}, conditional at the
explicit positive-real and moment scopes stated there. It has not
bounded the labelled completed C8.6 kernels in their required complex
activity norm. In particular, the source contains the prescribed
$|Z|^{-1}$, and the contact factor $2^{m-1/2}$ is retained. No
uniform control of either factor, of the inherited deficit, or of
the score-family cost follows from this calculation.

\subsection{The Full Quotient Response and its Complex Source Neighbourhood}
\label{sec:c8-fourth-full-response}

We now bound the full fourth derivative in
\eqref{eq:c8-fourth-defined-test}, including its denominator terms.
Keep the actual positive real bracket $\nu$, unlabelled inherited
insertion $b=\bar b_J(H)$, and a fixed marked contact $r$ from
Section~\ref{sec:c8-fourth-deficit}. All inherited variables and
operations are unchanged. Put
\begin{equation}\label{eq:c8-fourth-source-size}
    C_J=\frac{2^{m-1/2}}{|Z|}\nu[\bar q_r]^{1/2}\;.
\end{equation}
Equation \eqref{eq:c8-fourth-inherited-deficit} gives
$\nu[|b|^2]^{1/2}\leq C_J$, uniformly in independent auxiliary
holonomies. The joint two-loop statement uses the same joint lift,
$m=m_1+m_2$ and $Z=Z_1Z_2$.

Let $x_1,\ldots,x_4$ be four actual bounded real scalar tests, with
$M_j=\|x_j\|_{L^\infty(\nu)}$. For $s\in\mathbb C^4$ define
\begin{equation}\label{eq:c8-fourth-tilt}
    d(s)=\nu[e^{s\cdot x}],\qquad
    F(s;H)=\frac{\nu[b(H)e^{s\cdot x}]}{d(s)},\qquad
    s\cdot x=\sum_{j=1}^{4}s_jx_j\;.
\end{equation}
Thus $D(s)=D_\eta d(s)$ is the actual deformed vacuum denominator;
no sourced expectation is divided out. Only the four boundary-source
coefficients vary. The background, old operations and auxiliary
holonomies are fixed.

Fix $0\leq\tau<\log2$ and suppose $\sum_j|s_j|M_j\leq\tau$.
Since $|e^{s\cdot x}-1|\leq e^\tau-1$ pointwise,
\begin{equation}\label{eq:c8-fourth-vacuum-lower}
    |d(s)|\geq2-e^\tau>0,\qquad
    |D(s)|\geq(2-e^\tau)D_\eta\;.
\end{equation}
Bounded tests and integrability of $b$ justify all parameter derivatives
by dominated convergence on compact parameter sets. The numerator and
denominator are entire, and their quotient is holomorphic on this
neighbourhood. If $M_j=0$, that test is zero almost surely and the
corresponding derivative vanishes. No division by $M_j$ is required.

Set
\begin{equation}\label{eq:c8-fourth-tilted-functional}
    \nu_s[f]=\frac{\nu[fe^{s\cdot x}]}{d(s)},\qquad
    K_\tau=\frac{e^\tau}{2-e^\tau}\;.
\end{equation}
The functional $\nu_s$ is generally complex. For every integrable $f$,
the original positive law and \eqref{eq:c8-fourth-vacuum-lower} give
\begin{equation}\label{eq:c8-fourth-tilt-domination}
    |\nu_s[f]|\leq K_\tau\nu[|f|],\qquad
    |F(s;H)|\leq K_\tau C_J\;.
\end{equation}
The latter uses $\nu[|b|]\leq\nu[|b|^2]^{1/2}$. No independence
of the source and the boundary factors is assumed.

Let $I_5=\{0,1,2,3,4\}$, $z_0=b$ and $z_j=x_j$ for $j>0$.
For the set partitions $\Pi(I_5)$ the exact identity is
\begin{equation}\label{eq:c8-fourth-cumulant}
    \partial_1\partial_2\partial_3\partial_4F(s;H)
        =\sum_{\pi\in\Pi(I_5)}(-1)^{|\pi|-1}(|\pi|-1)!
            \prod_{A\in\pi}\nu_s\left[\prod_{j\in A}z_j\right]\;.
\end{equation}
To prove this finite cumulant formula without a sourced normaliser,
write $F(s+u)=\nu_s[be^{u\cdot x}]/\nu_s[e^{u\cdot x}]$ and expand
the reciprocal denominator as the geometric series in
$\nu_s[e^{u\cdot x}]-1$ near zero. The derivative labels acting on
the numerator join label $0$ in one block. The remaining labels form
nonempty denominator blocks. For $k-1$ such blocks their order in the
geometric series gives $(k-1)!$ copies with sign $(-1)^{k-1}$.
These are exactly the partitions in \eqref{eq:c8-fourth-cumulant}.
All terms of \eqref{eq:c8-fourth-quotient}, including the fourteen
proper-subset denominator terms, are retained.

There is exactly one source-containing block. By
\eqref{eq:c8-fourth-tilt-domination}, its moment is bounded by
$K_\tau C_J$ times its product of $M_j$; every other block costs
$K_\tau$ times its product of $M_j$. The numbers of partitions of
five labels into $k=1,\ldots,5$ blocks are $1,15,25,10,1$.
Multiplying these numbers by $(k-1)!$ therefore proves
\begin{equation}\label{eq:c8-fourth-complex-response-bound}
    |\partial_1\partial_2\partial_3\partial_4F(s;H)|
        \leq C_J\left(\prod_{j=1}^{4}M_j\right)
            \left(K_\tau+15K_\tau^2+50K_\tau^3
                +60K_\tau^4+24K_\tau^5\right)\;.
\end{equation}
At zero source, $K_0=1$ and this polynomial is $150$. With
$x=(p_i,q_i,p_j,q_j)$, summing the scalar bounds gives
\begin{equation}\label{eq:c8-fourth-full-family-bound}
    \left|\sum_{i,j}F^{ij}_{[4]}(0;H)\right|
        \leq150\,\frac{2^{m-1/2}}{|Z|}\mathcal M^2
            \nu[\bar q_r]^{1/2}\;.
\end{equation}
This controls the full quotient derivative without adding an unestimated
$\mathcal D_{\bar b_J}$ to it. The triangle inequality explicitly
pays the absolute family cost $\mathcal M^2$; no covariance-matrix or
tree-decay gain is inferred.

The same proof applies directly to the full unexpanded positive physical
law $\nu_{\mathrm{phys}}=\mathcal I_{\mathrm{phys}}/D_h$ if the
source has its stated positive chronological lift and the four tests
have finite essential suprema on that whole physical space. The literal
physical numerator and denominator are then
\begin{equation}\label{eq:c8-fourth-full-physical-tilt}
    N_{\mathrm{phys}}(s)=\mathcal I_{\mathrm{phys}}[be^{s\cdot x}],\qquad
    D_{\mathrm{phys}}(s)=\mathcal I_{\mathrm{phys}}[e^{s\cdot x}]\;,
\end{equation}
and $|D_{\mathrm{phys}}(s)|\geq(2-e^\tau)D_h$.
Every history, Haar factor, cutoff, inherited action, replacement
prefactor and postponed integration stays in $\mathcal I_{\mathrm{phys}}$.
Physical holonomies are substituted before integrations reading them.
Fixed-data score bounds do not imply uniform bounds on this physical
space, so this is an explicit conditional application, not a proof of
the missing history uniformities.

For the first connected caller WC-C7.5, one bounded boundary test $W$
gives $G(s)=\nu[be^{sW}]/\nu[e^{sW}]$ and
$G'(0)=\operatorname{Cov}_\nu(b,W)$. The same argument gives a joint
complex-source neighbourhood for $|s|\|W\|_\infty<\log2$, with the
actual vacuum denominator and $|G(s)|\leq K_\tau C_J$ when
$|s|\|W\|_\infty\leq\tau$. It uses a bound for the actual source
integrand through $C_J$, not merely its integrated mean. Its size and
localisation remain to be controlled at the required uniform scope.
This application is coefficientwise. Applying it to the complete source
$S$ in WC-C7.5 requires accounting for all of its marked and unmarked
contributions in their actual joint integration. No summability estimate
for that complete source is supplied by the bound for one $b_J$.

Complex boundary-source parameters and complex background variables
are distinct domains. This calculation does not bound the labelled
complex-background activities of WC-C8.6, nonlinear field-rotation
derivatives or cutoff-domain variations. It still pays $|Z|^{-1}$,
$2^{m-1/2}$ and the absolute family cost, and supplies no new ultraviolet
decay or noncommuting YM change-of-variables identity.

\subsection{Four Boundary Marks in the Labelled C8 Expansion}
\label{sec:c8-fourth-labelled}

The remaining caller is the sum of the completed kernels WC-C8.6,
with their inherited labels, rather than one unlabelled conditional
expectation. The calculation in this subsection differentiates those
very kernels. It identifies the cost of the four marks; it does not
assume a uniform bound for their inherited source seeds.

\paragraph{Comparison with the newer first-step records.}
The records \path{WILSON_IT1_LINE_DEFECT_2026_10_21.md},
Sections 3d--3e, and
\path{WILSON_IT1_D2_STEP_MAP_2026_10_22.md},
D2.0--D2.1, contain proposed upgrades of V3. The first extends the
closed-contour first-step representation to uncovered branches and
derives a complex-background extension with revised constants. Its
Section 3e(3.8) retains exponential growth in the length of the
undifferentiated contour words; only a group's own word cost is paid
by its local reserve. The second still assigns the old-insertion map
and the later R-operation to D2.2 and D2.7. Consequently these records
do not supply the uniform bound for the inherited labelled
$\bar b_{J,\sigma}$ required here. This is a comparison of operands
and scopes, not a rejection of their first-step derivations. No open
matrix insertion from those records is adopted as an inductive carrier.

\paragraph{Actual factors and four independent marks.}
Fix the outer data, $J$, $\sigma$ and $H$, and keep
$b=\bar b_{J,\sigma}(H)$ inside the component operations. For each
input boundary label $a$, write $u_a=\operatorname{Re}\beta_a$ on
the real integration variables. Let $B_a$ be a finite uniform bound
for $|\beta_a|$ on all the integration variables of that factor,
at these fixed outer data. This paragraph assumes $0\leq B_a\leq1$.
The small input boundary bounds used in LFR2 (1.103) are of this
form. No such bound on the source itself is inferred from them.
Choose four labels $a_1,\ldots,a_4$, with repetitions allowed, and set
\begin{equation}\label{eq:c8-fourth-marked-edge}
    v_a(s)=\exp\left(\beta_a+
        \sum_{r=1}^{4}s_r\mathbf 1_{a=a_r}u_a\right)-1\;.
\end{equation}
Then $\prod_a(1+v_a(s))=\exp(B_\eta+\sum_rs_ru_{a_r})$.
Thus the numerator and vacuum denominator are exactly those in
\eqref{eq:c8-fourth-defined-test}. The source, the component
functionals and all earlier denominator labels stay fixed during
this current boundary-source differentiation. For a nonempty subset
$A\subseteq\{1,2,3,4\}$, differentiation gives
\begin{equation}\label{eq:c8-fourth-edge-derivative}
    \partial_Av_a(0)=e^{\beta_a}u_a^{|A|}
        \prod_{r\in A}\mathbf 1_{a=a_r}\;.
\end{equation}
Put $w_a=B_ae^{B_a}$. Since $|u_a|\leq B_a\leq1$, the
undifferentiated factor and every nonzero mixed derivative through
order four have modulus at most $w_a$. When $B_a=0$, both the factor
and these derivatives vanish; no division by $B_a$ is used. This is
analyticity in $s$ only. Taking a real part does not furnish a
holomorphic extension in the background.

\paragraph{Terms, multiplicities and their majorants.}
Resolve WC-C8.6 into the following terms, without altering its labels.
A term $\mathcal T$ specifies the source component $Z_0$ and its
connected selected boundary family $D_0$, an ordered list of $q$
vacuum clusters, the ordered polymer vertices in each cluster, and
the connected selected boundary family $D_{rj}$ at each such vertex.
The source family is allowed to be empty. The vacuum families are
nonempty. Each vacuum cluster touches $Z_0$, and the union of all
their supports with $Z_0$ is the assigned support $X$.
All these supports contain every whole inherited $Y_i$ that their
integrands read, as prescribed in WC-C8.6.

For a cluster with ordered vertices $(Z_1,\ldots,Z_n)$, let $G$
range over connected graphs on $\{1,\ldots,n\}$ and retain
the graph sum
\begin{equation}\label{eq:c8-fourth-ursell}
    U(Z_1,\ldots,Z_n)=
        \sum_G\prod_{\{j,l\}\in E(G)}
            \left(-\mathbf 1_{Z_j\not\sim Z_l}\right)\;.
\end{equation}
It is this entire sum, including repeated labels and
self-incompatibility, whose modulus is taken below. Taking the
modulus of each graph before summing would give a different and
potentially divergent majorant.

Let $I(D)$ be the component slots read by a selected family $D$;
let $I_0$ include also every slot read by $b$. Write $\nu_i$ for the
literal normalised component measure of WC-C8.4, $|\nu_i|$ for its
total variation, and $K_i=|\nu_i|(1)$. These are the original
integral functionals, with the original normalising scalars.
LFR2 (1.73)--(1.75) gives $K_i\leq e^3$ on its stated domain.
Define nonnegative, possibly infinite, quantities
\begin{equation}\label{eq:c8-fourth-selected-majorants}
    \begin{aligned}
        M_b(D_0)&=
            \left(\prod_{i\in I_0}|\nu_i|\right)[|b|]
                \prod_{a\in D_0}w_a,\\
        M_0(D)&=\left(\prod_{i\in I(D)}K_i\right)
                \prod_{a\in D}w_a\;.
    \end{aligned}
\end{equation}
Every numerator, cutoff and earlier scalar denominator in $b$
remains inside the first integral. In particular $M_b(D_0)$ is not
replaced by the modulus of an already integrated source coefficient.
Taking total variation defines a majorant, not a new physical law.

If cluster $r$ has $n_r$ vertices, set
\begin{equation}\label{eq:c8-fourth-term-majorant}
    A_{\mathcal T}=
        \frac{1}{q!}\prod_{r=1}^{q}
            \frac{|U(Z_{r1},\ldots,Z_{rn_r})|}{n_r!}
        M_b(D_0)\prod_{r=1}^{q}\prod_{j=1}^{n_r}M_0(D_{rj})\;.
\end{equation}
Incompatible pairs in $U$ have precisely the original compatibility
rule. The sign $(-1)^q$ of denominator completion is retained in
the actual term and removed only in $A_{\mathcal T}$.
Repeated occurrences of a boundary label in different vertex
integrals are distinct occurrences. Let $N_a(\mathcal T)$ count
them, including the source vertex, and put
$n(\mathcal T)=\sum_aN_a(\mathcal T)$.

For each assignment of the four derivative labels to factor
occurrences, apply \eqref{eq:c8-fourth-edge-derivative} inside
the original integrals. The product rule sums over all such
assignments, including several derivatives on one occurrence.
It gives the termwise estimate
\begin{equation}\label{eq:c8-fourth-term-derivative}
    |\partial_1\partial_2\partial_3\partial_4
        \mathcal T(0)|
        \leq A_{\mathcal T}
            \prod_{r=1}^{4}N_{a_r}(\mathcal T)\;.
\end{equation}
The integrations factor only across the component slots already
disjoint in the C8 construction. No probabilistic independence
inside an inherited component is introduced. Every differentiated
integrand is a closed scalar contraction multiplied by invariant
scalar boundary factors; invariance is preserved before the estimate.

\paragraph{The test family and the occurrence reserve.}
For the same finite family $(\alpha_i,\gamma_i)$ as above, define
$h$ to be the largest multiplicity of an ordered pair $(a,c)$ in
that list, and set $h=0$ for the empty family. Then
\begin{equation}\label{eq:c8-fourth-family-count}
    \begin{aligned}
        \sum_{i,j}\prod_{a\in(\alpha_i,\gamma_i,\alpha_j,\gamma_j)}
            N_a(\mathcal T)
            &=\left(\sum_iN_{\alpha_i}(\mathcal T)
                N_{\gamma_i}(\mathcal T)\right)^2\\
            &\leq h^2 n(\mathcal T)^4\;.
    \end{aligned}
\end{equation}
The product on the first line is a product over four positions,
with repeated labels retained. The inequality follows by replacing
each ordered-pair multiplicity by $h$ and summing all ordered pairs
of boundary labels. It does not multiply by the total number of
tests in the ambient volume.

For $\delta>0$, let
\begin{equation}\label{eq:c8-fourth-boosted-majorant}
    \mathcal A^{\delta}_{J,\sigma,X}
        =\sum_{\mathcal T:\,\operatorname{supp}\mathcal T=X}
            A_{\mathcal T}e^{\delta n(\mathcal T)}\;.
\end{equation}
Equivalently this is the same absolute C8 majorant with $w_a$
replaced by $e^\delta w_a$ at each occurrence. It does not change
the measure, the source normalisation, the physical denominator or
the support. Since $t^4e^{-\delta t}\leq(4/(e\delta))^4$ for
$t\geq0$, \eqref{eq:c8-fourth-term-derivative} proves
\begin{equation}\label{eq:c8-fourth-labelled-bound}
    \sum_{i,j}\sum_{\mathcal T:\,\operatorname{supp}\mathcal T=X}
        |\partial_{[4]}\mathcal T^{ij}(0)|
        \leq h^2\left(\frac{4}{e\delta}\right)^4
            \mathcal A^{\delta}_{J,\sigma,X}\;.
\end{equation}
This is an inequality of nonnegative sums, allowing infinity.
Whenever its right-hand side is summable, it controls the full
labelled derivative, including all touching-denominator terms.
The empty source family with no vacuum clusters has zero derivative;
it contributes no exception at $n=0$.

Here is the convergence justification without a formal exchange of
infinite sums. For $\max_r|s_r|\leq\delta/16$, each deformed factor
satisfies
\begin{equation}\label{eq:c8-fourth-source-polydisc}
    |v_a(s)|\leq
        \left(1+\frac{\delta}{4}\right)
        e^{\delta/4}w_a\leq e^{\delta/2}w_a\;.
\end{equation}
Thus finiteness of the summed $\mathcal A^\delta$, with the vacuum
cluster expansion valid for these enlarged activities, gives local
uniform absolute convergence on this polydisc. The same occurrence
estimate with the unused half of the reserve controls differentiated
series on smaller closed polydiscs. Termwise differentiation is then
valid. The vacuum cluster identity represents the original
$D_\eta(s)$ there; its exponential is nonzero. No source-containing
partition function is used for normalisation.

The estimate remains valid after multiplication by any specified
nonnegative support weight and summation over $X,J,\sigma$, with
the same constant. Consequently the derivative cost alone can be
paid by a small strengthening of the boundary activity parameter.
In the LFR2 (1.103) estimate, the parameter $O(1)c_1$ must be chosen
below its convergence threshold also after this strengthening.
Neither $e^3$ per inherited component nor its size is raised to a
new fourth-order power by this calculation. This uses an occurrence
reserve, not a claimed tree decay inside $Y_i$: p.~391 explicitly
excludes that decay from the vacuum argument because the $Y_i$ are
fixed there.

\paragraph{Literal caller and the remaining estimate.}
Write
$R_\eta=\sum_{J,\sigma,X,i,j}\partial_{[4]}
f^{ij}_{J,\sigma,X}(0;H^{\rm phys})$ whenever that sum exists,
with physical holonomies substituted before any integration that
reads them. In this current bracket the transports defining
$H^{\rm phys}$ use the frozen $V$, as stated after WC-C8.3;
they are therefore legitimate fixed arguments of the current
component integrals. Their dependence is restored before subsequent
RG integrations. The following bound uses a source majorant at
these same fixed arguments, not a supremum taken after integrating
a variable holonomy. Its physical insertion is
\begin{equation}\label{eq:c8-fourth-labelled-physical}
    \frac{\mathcal O_\eta[A_\eta D_\eta R_\eta]}{D_h}\;.
\end{equation}
This is a conditional-response insertion in the physical integral.
It is not asserted to equal the fourth derivative of the full
physical quotient: derivatives of the outer vacuum weights would
also be required for that statement.
If $\mathcal O_\eta^{\rm abs}$ denotes total variation of the
literal outer sums and integrations, the bound obtained by
substitution is
\begin{equation}\label{eq:c8-fourth-labelled-physical-bound}
    \frac{|\mathcal O_\eta[A_\eta D_\eta R_\eta]|}{D_h}
        \leq \frac{h^2}{D_h}
            \left(\frac{4}{e\delta}\right)^4
            \mathcal O_\eta^{\rm abs}\left[
                |A_\eta D_\eta|
                \sum_{J,\sigma,X}\mathcal A^\delta_{J,\sigma,X}
            \right]\;.
\end{equation}
The source renormalisation $Z$ is still in every source seed. The
right-hand side requires its own finite, uniform physical estimate;
the factor $D_h$ has not disappeared. In particular the source-free
LFR2 bounds do not estimate
$(\prod_{i\in I_0}|\nu_i|)[|\bar b_{J,\sigma}|]$, its sum over
earlier denominator labels, or its complex-background dependence.
The new calculation pays the four current boundary marks locally
once the actual source majorant is available. It neither supplies
that missing majorant nor creates an ultraviolet small factor.

The immediate falsification checks are algebraic. With no selected
boundary occurrence, the derivative is zero. Four derivatives can
hit a single occurrence, so four distinct boundary labels must not
be required. Repeated labels in different denominator vertices
must remain separate occurrences. If the closed source integrand
vanishes pointwise, every differentiated term vanishes as well;
flatness of the background alone is not identified with this
pointwise condition after fluctuation integration.

\subsection{A First-Step Source Block with the Coarse Comparison Retained}
\label{sec:c8-fourth-first-seed}

This calculation addresses the source factor left unevaluated in
\eqref{eq:c8-fourth-selected-majorants}. Its scope is a first-step
selected block, with the actual retained data, followed by multiplication
by a later boundary test that reads those retained data. It is not
an estimate for each separate C8 coefficient after absolute values.
The finite sum over the marked choices in one interval is taken
inside the closed source contraction. No measure is changed.

\paragraph{First-step input and its scope.}
Fix real retained data $y$ and a nonempty selected group $W$ in V3,
Lemma \texttt{c5-5-3}, at $t=1$. Use the same selected integration
operation, denoted here by $\mathscr Q_W$, that defines
$c_W(h)=c_W^1(h)$ in Lemmas \texttt{c5-5-7} and
\texttt{c5-5R-3}. In particular, if $\mathcal A(W)$ is its root
family, $\ell_R$ are the selected labels of root $R$, and
$\mathsf J_R\neq\varnothing$ their free-bond probe sets, this operation is
\begin{equation}\label{eq:c8-fourth-seed-operation}
    \mathscr Q_W[g]
        =w_\Lambda^1(W)
            \sum_{(\ell_R)_{R\in\mathcal A(W)}}
                \int \Delta_{\mathsf J}g\,
                    \dd\bigotimes_{R\in\mathcal A(W)}
                        \Lambda_{\Lambda;\ell_R}^{f,1}\;,
    \qquad
    \mathsf J=\bigsqcup_{R\in\mathcal A(W)}\mathsf J_R\;.
\end{equation}
Here $\Delta_{\mathsf J}$ is the finite difference in the young
free-bond variables about their real base point $m$. The measures,
translated cutoffs, exterior shift, selected host factors and
touching vacuum-cluster weight $w_\Lambda^1(W)$ are V3's original
ones. This is a signed linear integration operation, not a
probability. Its nonempty probe sets imply
$\mathscr Q_W[1]=0$.

Let $T_W(B)$ be the fine ordered transport along the contact interval
$\mathcal H(W)$ and $P_W=T_W(0)$. The source input and estimates used
below are, for every $h\in G$,
\begin{equation}\label{eq:c8-fourth-seed-input}
    \begin{aligned}
        c_W(h)&=\mathscr Q_W[
            \chi(T_W(B)P_W^{-1}h)],\\
        |c_W(h)|&\leq M_W^{\rm ph},\qquad
        |a_W|\leq M_W^{\rm ph},\qquad
        d_W=\sup_{h\in G}|c_W(h)-a_W\chi(h)|\;.
    \end{aligned}
\end{equation}
The first two bounds are precisely the real, pointwise-in-$y$
input of V3 Lemma \texttt{c5-5R-3}; in particular
$d_W\leq2M_W^{\rm ph}$. They require that lemma's actual hypotheses,
constants, cutoffs and small-coupling conditions. Their proof uses
signed source differences and parity before total variation, and
cannot be replaced by a bound on the raw source's absolute integral.
The newer IT-1 Section 3d proposes their extension to uncovered
branches. The calculation below applies wherever the stated
first-step input holds; it does not silently assume an all-history
extension or use a failed-coverage probability to discard other
branches. The corresponding stronger deviation estimate for $d_W$
is a separate input in IT-1 Section 3a.

\paragraph{The C8 polynomial inside the same selected operation.}
Choose the smallest coarse-bond interval $I$ containing
$\mathcal H(W)$, with at most $L-1$ added fine traversals at either
end. Here the group is in a genuine straight or corner germ; other
groups retain their original multi-holonomy form. Let $P_C$ be
the transport of the frozen coarse $V$ along $I$, and let
$F_a(B)$ be the fine transport along coarse edge $a$ in the
reconstructed original fine field, with the selected young
variables of $W$ active and all other young variables at $m$.
Write $C_a=V(c_a)$ and $H\in G$ for a closed holonomy at the
initial point of $I$. The closed C8 polynomial is
\begin{equation}\label{eq:c8-fourth-local-polynomial}
    \mathcal P_I(z,B;H)
        =\frac{1}{Zd}\operatorname{Tr}\left[
            \left(\prod_{a\in I}
                (R(C_a)+z_a(R(F_a(B))-R(C_a)))\right)
            R(P_C)^{-1}R(H)\right]\;.
\end{equation}
The product has the contour order. Its coefficient
$\mathcal P_{I,J}$ is the C8 coefficient with marked set $J$
inside $I$, after the internal coarse holonomies and tree
transports are restored as in WC-C8.7. This equality follows by
telescoping those tree transports inside the whole trace.
All coefficients are closed invariant contractions; none of the
open factors in this algebraic polynomial is assigned an activity
norm.

The empty coefficient is $Z^{-1}\chi(H)$ and is independent of
the young variables. Let $T_I(B)=\prod_{a\in I}F_a(B)$ and
$P_I=T_I(0)$. Finite summation and the linear operation give
\begin{equation}\label{eq:c8-fourth-local-resum}
    \sum_{\varnothing\neq J\subseteq I}
        \mathscr Q_W[\mathcal P_{I,J}(\,\cdot\,;H)]
        =\frac{1}{Z}\mathscr Q_W[
            \chi(T_I(B)P_C^{-1}H)]\;.
\end{equation}
The field on the right is the reconstructed fine field inside
the selected integral, not its minimiser. Any earlier finite
endpoint decomposition used to represent this same first-step
source is summed to its physical endpoint before this identity
is used. This does not bound the sum of the absolute values of
those endpoint coefficients.

By the definition of the contact set, selected young variables of
$W$ do not affect the added traversals. Thus
$T_I(B)=P_L T_W(B)P_R$ and $P_I=P_LP_WP_R$, with $P_L,P_R$
fixed during $\mathscr Q_W$. Put $K_I=P_IP_C^{-1}$, a closed
relative holonomy. Cyclicity of the trace proves
\begin{equation}\label{eq:c8-fourth-seed-identification}
    \sum_{\varnothing\neq J\subseteq I}
        \mathscr Q_W[\mathcal P_{I,J}(\,\cdot\,;H)]
        =\frac{1}{Z}c_W(P_L^{-1}K_IHP_L)\;.
\end{equation}
The argument of $c_W$ is a closed loop based at the initial
point of $\mathcal H(W)$. Under the induced gauge transformation
it is conjugated there. This identity uses the unitary $P_L$
only to change the base point of a complete closed contraction.

\paragraph{The scalar part and a pointwise source estimate.}
Retain the existing reference coefficient $a_W/Z$ and define
the resummed subtracted source block by
\begin{equation}\label{eq:c8-fourth-resummed-seed}
    G_W(H)=
        \sum_{\varnothing\neq J\subseteq I}
            \mathscr Q_W[\mathcal P_{I,J}(\,\cdot\,;H)]
        -\frac{a_W}{Z}\chi(H)\;.
\end{equation}
This scalar extraction is inside the selected marked expansion.
It changes neither the prescribed $Z$ nor a perimeter or corner
coefficient. There is only one global factor $Z^{-1}$ for the
Wilson insertion. Formula \eqref{eq:c8-fourth-resummed-seed}
describes a single-group contribution; multiplying several such
scaled scalar kernels would insert spurious powers of $Z^{-1}$
and would also lose their joint closed contraction. No such
product or family reconstruction is asserted here.
Character invariance under conjugation gives the
exact decomposition
\begin{equation}\label{eq:c8-fourth-seed-coarse-defect}
    ZG_W(H)=
        c_W(P_L^{-1}K_IHP_L)-a_W\chi(K_IH)
        +a_W(\chi(K_IH)-\chi(H))\;.
\end{equation}
In particular the blocked coarse word cannot be replaced by $P_I$.
For $q(K)=1-\operatorname{Re}\chi(K)$, the closed-trace
Cauchy--Schwarz estimate gives
$|\chi(KH)-\chi(H)|\leq\sqrt{2q(K)}\leq2$ for $K,H\in G$.
Consequently, at each retained real datum,
\begin{equation}\label{eq:c8-fourth-seed-pointwise}
    \sup_{H\in G}|G_W(H)|
        \leq\frac{d_W+|a_W|\sqrt{2q(K_I)}}{|Z|}
        \leq\frac{4M_W^{\rm ph}}{|Z|}\;.
\end{equation}
There is no factor $2^{|I|}$: the local marked choices were
summed exactly before estimation. At the consistent flat reference
data, $K_I=1$ and $c_W(h)=a_W\chi(h)$, so the residual vanishes.
A nonzero exterior shift is retained in $P_I$ and is not declared
flat merely because the unshifted background is flat.

\paragraph{A later score does not undo an integration already performed.}
Let $\mathcal L$ be an actual later inherited integration on the
remaining variables $y$, and let $T(y)$ be its boundary test.
Let $\mathscr E$ be the measurable part of those data on which the
first-step input \eqref{eq:c8-fourth-seed-input} holds, with the
same constants. All its original indicators are retained.
The young variables integrated by $\mathscr Q_W$ have already been
eliminated; they are not given a fresh probability and do not become
arguments of $T$. Inequality \eqref{eq:c8-fourth-seed-pointwise}
is pointwise in $y$, uniformly in $H$. Whenever the displayed
majorant is integrable, it implies
\begin{equation}\label{eq:c8-fourth-seed-score}
    |\mathcal L[\mathbf 1_{\mathscr E}G_W(H(y))T]|
        \leq\frac{1}{|Z|}|\mathcal L|\left[
            \mathbf 1_{\mathscr E}
            (d_W+|a_W|\sqrt{2q(K_I)})|T|\right]\;.
\end{equation}
Here $|\mathcal L|$ is the total variation of the original
integration operation, with all its original scalar normalisers.
The possibly field-dependent $H(y)$ is evaluated before this later
integration. The uniform bound is on the integrand at each $y$,
so no supremum is moved through an integral. This is a quantitative
source--score bound for these first-step resummed blocks. It uses
the stronger pointwise statement of V3, rather than an integrated
mean alone. It is not a proof that a current unintegrated source
has the parity of an earlier eliminated young variable.
For the unrestricted physical integral the exact complementary
term $\mathcal L[\mathbf 1_{\mathscr E^c}G_W(H(y))T]$ remains.
It is not bounded by a small probability or omitted. Using the
newer uncovered-branch extension would require verifying its
scope at these same retained data.

For example, the source-root factor of
\eqref{eq:c8-fourth-selected-majorants}, with $b$ replaced by this
actual regrouped block, satisfies
\begin{equation}\label{eq:c8-fourth-seed-in-majorant}
    M_{\mathbf 1_{\mathscr E}G_W}(D_0)
        \leq\frac{1}{|Z|}
            \left(\prod_{i\in I_0}|\nu_i|\right)
                [\mathbf 1_{\mathscr E}
                    (d_W+|a_W|\sqrt{2q(K_I)})]
            \prod_{a\in D_0}w_a\;.
\end{equation}
The boundary factors remain inside the same inherited operations
until this pointwise majorisation is applied. Thus the four-mark
estimate can be used for these regrouped seeds while retaining
their first-step source gain and the physical denominator in
\eqref{eq:c8-fourth-labelled-physical-bound}.

\paragraph{Estimating the retained coarse-word defect on regular rows.}
This paragraph uses the averaging formula RG1 (0.4), the choice
made on its chart domain in V3. Let $U_*$ be the real fine field
at the base point of the selected integral, so that
$P_I=U_*(I)$. On the rows considered assume the actual constraint
$M(U_*)(c_a)=V(c_a)$ and an injective coordinate box $Q_a$
containing the two adjacent blocks and the paths in RG1 (0.4).
For a negatively oriented traversal, use the inverse of the
positive-bond constraint; its relative error is an inverse
conjugate and has the same norm.
Write $D=4$ for the lattice dimension, to distinguish it from
the representation dimension $d$. Define
$\delta_a=\max_{p\subset Q_a}\|U_*(\partial p)-1\|_{\rm op}$
in the fixed faithful matrix realisation of $G$.
Let $r_*>0$ be at most $1/2$ and small enough that the matrix
ball of radius $r_*$ in $G$ lies in the stipulated logarithm
chart. Suppose $D^2L^2\delta_a\leq r_*$.

For a point $x$ in the initial block and paths $\Gamma,\Gamma'$
of RG1 p.~252, put
$L_{a,x,\Gamma,\Gamma'}=
U_*(\Gamma\cup[x,x']\cup(-\Gamma')\cup(-c_a))$.
These are complete loops based at the initial point of $c_a$.
Each has a nonabelian plaquette filling using at most $D^2L^2$
plaquettes, counted with multiplicity. To see the count, put each
of the two shortest coordinate paths into one fixed coordinate
order by interchanging adjacent coordinate runs. This requires
at most $D(D-1)/2$ rectangular fillings per path, each of area at
most $L^2$. The two equally ordered paths are then joined by
at most $D-1$ rectangles swept along the coarse bond, also of
area at most $L^2$. The total is at most
$(D^2-1)L^2$. These rectangles lie in $Q_a$. The exact ordered
plaquette product has unitary conjugating transports, hence
\begin{equation}\label{eq:c8-fourth-average-loop-bound}
    \|L_{a,x,\Gamma,\Gamma'}-1\|_{\rm op}
        \leq D^2L^2\delta_a\;.
\end{equation}
This is a bound on complete loop holonomies, before an average
or an absolute source estimate.

Let $\omega_{x,\Gamma,\Gamma'}$ be the nonnegative weights
$L^{-D}|\mathbf G(c_-,x)|^{-1}
|\mathbf G(c_+,x')|^{-1}$ of RG1 (0.4); their sum is one.
That formula and the actual constraint give
\begin{equation}\label{eq:c8-fourth-average-relative}
    E_a=\sum_{x,\Gamma,\Gamma'}\omega_{x,\Gamma,\Gamma'}
            \log L_{a,x,\Gamma,\Gamma'},\qquad
    V(c_a)=e^{E_a}U_*(c_a),\qquad
    U_*(c_a)V(c_a)^{-1}=e^{-E_a}\;.
\end{equation}
Here $E_a$ is skew-Hermitian and transforms by conjugation at
the row base. Let $C_R$ be the norm of the differential of $R$
from the fixed Lie-algebra operator norm to the representation
operator norm. The convergent logarithm series gives
$\|\log L_{a,x,\Gamma,\Gamma'}\|_{\rm op}
\leq2\|L_{a,x,\Gamma,\Gamma'}-1\|_{\rm op}$.
Integrating the derivative of the unitary exponential gives
$\|R(e^{-E_a})-1\|_{\rm op}\leq C_R\|E_a\|_{\rm op}$.
The product over rows expresses $K_I$ as the ordered product
of these closed relative holonomies, conjugated by coarse
prefix transports. Telescoping that closed product yields
\begin{equation}\label{eq:c8-fourth-coarse-curvature}
    \sqrt{2q(K_I)}
        \leq \|R(K_I)-1\|_{\rm op}
        \leq \min\left(2,\,
            2C_RD^2L^2\sum_{a\in I}\delta_a\right)\;.
\end{equation}
The first inequality follows from
$2q(K_I)=d^{-1}\|R(K_I)-1\|_{\rm HS}^2$.
In particular substitution into
\eqref{eq:c8-fourth-seed-pointwise} gives a curvature estimate
for the actual retained coarse-word term, with no replacement
of $P_C$ by $P_I$. Rows not satisfying these regularity and
geometry conditions retain the original $\sqrt{2q(K_I)}\leq2$
bound; they are not discarded.

This computation does not identify $U_*$ with the unshifted
minimiser. Its curvature includes the retained exterior shift.
Turning the displayed row-curvature sum into the uniform
local deviation and attached-part bounds required by the
induction still needs the corresponding mean-field and
support estimates. It also does not extend the real logarithm
argument to complex backgrounds.

This substitution does not bound the earlier, stronger direct sum
of absolute values over separate $J,\sigma$. It supplies a bound
after the specified finite C8 regrouping. To obtain an induction
one still needs the simultaneous family regrouping, the complete
supports in coarse variables, the complex-background extension,
and the next-step source estimate. The relative holonomy $K_I$
also needs its curvature and exterior-response estimate at that
scope. None of these requirements is removed by the bound
$4M_W^{\rm ph}/|Z|$.

\subsection{Simultaneous First-Step Families on the Coarse Contour}
\label{sec:c8-fourth-coarse-families}

The single-group bound cannot be multiplied to reconstruct the
source. This subsection instead starts from V3's actual simultaneous
family identity and keeps its joint reference exclusion factor.
It is a further conditional first-step calculation on the real
retained data for which those inputs hold. The physical contribution
from other data is retained as in
\eqref{eq:c8-fourth-seed-score}.

\paragraph{Closed kernels and the actual family input.}
Use the already defined algebra $E_\chi$, its product $\star$,
and its norm $\|\cdot\|_\star$ from
\path{WILSON_RG_COMPARISON.tex}, equations
\texttt{eq:Ecoefficient} and \texttt{eq:closedstar}.
The norm is the operator
norm of left $\star$ multiplication on $E_\chi$ with its Haar
$L^2$ norm. The unit is $\chi$, closure is $\tau(f)=f(1)$,
the norm is submultiplicative, and conjugation by a real unitary
transport is isometric. These are identities on complete closed
functions. No open segment is made into an inductive activity.

At the base point of each fine group $W$, put
$r_W(h)=c_W(h)-a_W\chi(h)$. Let $U_*$ be the actual base fine
field, with its retained exterior data. V3
\texttt{eq:c5-5-5} and Lemma \texttt{c5-5-7}(b) give
\begin{equation}\label{eq:c8-fourth-input-families}
    F^\Lambda(1)=S^\circ_\Lambda\left[
        \chi(U_*(C))+
        \sum_{\mathfrak B\neq\varnothing}
            q^\circ_\Lambda(\mathfrak B)
            \mathcal R(\mathfrak B;U_*)\right]\;.
\end{equation}
Here $\mathfrak B$ ranges over the original compatible fine
groups. $\mathcal R$ is their ordered, complete closed contraction
with the fine base transports, not a product of characters.
Let $\operatorname{Hit}(\mathfrak B)$ be the reference clusters
incompatible with at least one original group in $\mathfrak B$.
The reference factor is exactly
\begin{equation}\label{eq:c8-fourth-joint-exclusion}
    q^\circ_\Lambda(\mathfrak B)
        =\exp\left(-\sum_{\Gamma\in\operatorname{Hit}(\mathfrak B)}
            \Phi^\circ_\Lambda(\Gamma)\right),\qquad
    |q^\circ_\Lambda(\mathfrak B)|
        \leq\prod_{W\in\mathfrak B}e^{\theta^\circ(W)}\;.
\end{equation}
A reference cluster hits the family if it is incompatible with
at least one original $W$. This definition is unchanged below.
The factor $S^\circ_\Lambda$, with its existing perimeter, corner
and carrier contributions, remains in the physical numerator.

Write $\sigma(W)\geq1$ for V3's size, $\nu>0$ for its reserve,
and let $\epsilon=\bar\varepsilon_{\rm ph}$. The specific
pointwise pin input needed here is, for every fine traversal $n$,
\begin{equation}\label{eq:c8-fourth-fine-pin-input}
    \sum_{W:\,n\in\mathcal H(W)}
        e^{2\nu\sigma(W)+\theta^\circ(W)}
            \|r_W\|_\star\leq\epsilon\;.
\end{equation}
This follows at the scope of V3 Lemma
\texttt{c5-5R-3} from
$\|r_W\|_\star\leq K_\chi d_W\leq2K_\chi M_W^{\rm ph}$
and its displayed pin bound. In particular
$\epsilon=O(g_0^2(1+r)^4)$ there, with V3's constants and
thresholds. This input is not asserted for a new RG scale.

\paragraph{The merging rule.}
All intervals in this paragraph are sets of traversal positions,
so visits of the same geometric bond would remain distinguished.
Let $I_W$ consist of the coarse traversal positions met by
$\mathcal H(W)$. For a fixed original compatible family
$\mathfrak B$, join $W$ and $W'$ precisely when
$I_W\cap I_{W'}\neq\varnothing$. Let $\mathcal D$ denote a
connected component of this overlap graph and put
$I_{\mathcal D}=\bigcup_{W\in\mathcal D}I_W$ and
$s(\mathcal D)=\sum_{W\in\mathcal D}\sigma(W)$.
The distinct $I_{\mathcal D}$ are disjoint coarse intervals.
If one is the entire contour, it is the only component.

This is an exact regrouping, not a change of compatibility.
A new family consists of overlap-connected sets $\mathcal D$
whose union is an original compatible family and whose coarse
intervals are disjoint. The union recovers $\mathfrak B$
uniquely. In particular original spatial incompatibilities are
retained even when two coarse intervals are disjoint.
The reference factor of this family is
$q^\circ_\Lambda(\bigcup_{\mathcal D}\mathcal D)$.
A reference cluster is not deleted merely because it intersects
an enlarged coarse interval; it is tested against the original
fine groups as in \eqref{eq:c8-fourth-joint-exclusion}.

\paragraph{The closed kernel of a proper merged interval.}
Suppose $I_{\mathcal D}$ is a proper interval. At its initial
coarse vertex $x_{\mathcal D}$, set
$P_{\mathcal D}=U_*(I_{\mathcal D})$,
$C_{\mathcal D}=V(I_{\mathcal D})$ and
$K_{\mathcal D}=P_{\mathcal D}C_{\mathcal D}^{-1}$.
For $W\in\mathcal D$, let $P_{\mathcal D,W}$ be the fine base
transport from $x_{\mathcal D}$ to the initial point of
$\mathcal H(W)$, inside $I_{\mathcal D}$, and set
\begin{equation}\label{eq:c8-fourth-merged-kernel}
    \begin{aligned}
        r_W^{\mathcal D}(H)
            &=r_W(P_{\mathcal D,W}^{-1}HP_{\mathcal D,W}),\\
        k_K(H)&=\chi(KH),\\
        g_{\mathcal D}
            &=\left(\mathop{\star}_{W\in\mathcal D}^{\rm ordered}
                    r_W^{\mathcal D}\right)\star k_{K_{\mathcal D}}\;.
    \end{aligned}
\end{equation}
The order is the order of the disjoint fine contact intervals
along $I_{\mathcal D}$. Each displayed kernel is a complete
closed function at its stated base. Its support includes all
original group footprints, the entire interval and every return
path used here. No localisation in the coarse field variables
is deduced just from the contour interval.

In a matrix realisation used only to verify this algebraic
identity, if $r_W=f_{A_W}$, the coefficient in
$g_{\mathcal D}$ is
\begin{equation}\label{eq:c8-fourth-merged-word}
    \left(\prod_{W\in\mathcal D}^{\rm ordered}
        R(P_{\mathcal D,W})A_WR(P_{\mathcal D,W})^{-1}\right)
        R(P_{\mathcal D})R(C_{\mathcal D})^{-1}\;.
\end{equation}
Multiplying by the coarse transport $R(C_{\mathcal D})$ in the
whole contour trace restores exactly the original fine base
word with all the insertions in $\mathcal D$ in their original
order. The product is not replaced by a product of its traces.
The already fixed closed-kernel norm therefore gives
\begin{equation}\label{eq:c8-fourth-merged-norm}
    \|g_{\mathcal D}\|_\star
        \leq\prod_{W\in\mathcal D}\|r_W\|_\star\;.
\end{equation}
The unitary factor $k_{K_{\mathcal D}}$ has norm one. If
$I_{\mathcal D}$ is the whole contour, keep the original
closed scalar $\mathcal R(\mathcal D;U_*)$ instead of cutting
an insertion at a coarse base point. Its modulus has the same
product bound. This case introduces no exterior holonomy.

\paragraph{The coarse edges outside the merged intervals.}
They cannot simply be replaced by coarse $V$ edges. For each
such coarse traversal $a$, retain
\begin{equation}\label{eq:c8-fourth-zero-root}
    K_a=U_*(c_a)V(c_a)^{-1},\qquad
    \ell_a(H)=\chi(K_aH)-\chi(H)\;.
\end{equation}
These are closed kernels and the exact difference identity is
$R(U_*(c_a))=(1+[R(K_a)-1])R(V(c_a))$.
Expanding this identity on the unused coarse edges sums over
subsets of those edges. On a selected edge insert $\ell_a$;
on the others keep the coarse transport. This is the sum over
the inactive C8 marked choices at the young base point, not
an additional integration or an approximation.
On the regular rows of
\eqref{eq:c8-fourth-coarse-curvature},
\begin{equation}\label{eq:c8-fourth-zero-root-bound}
    \|\ell_a\|_\star
        =\|R(K_a)-1\|_{\rm op}
        \leq\min(2,\,2C_RD^2L^2\delta_a)\;.
\end{equation}
Other rows retain the bound two. At consistent flat data these
kernels vanish. They carry no additional power of $Z^{-1}$.

For clarity, let $\mathcal C_V$ denote the ordered complete
closed contraction of the kernels $g_{\mathcal D}$ and
$\ell_a$ on their disjoint coarse intervals, with the
remaining coarse $V$ transports. It is defined by
\eqref{eq:c8-fourth-merged-word} and the edge difference above
inside a single normalised trace. Equivalently it is the same
contraction in the closed algebra $E_\chi$. Only this whole
contraction is estimated. For real $V$, its modulus is bounded
by the product of the closed-kernel norms, because closure has
norm at most one and the intervening unitary transports act
isometrically.

Let $\mathfrak D$ range over the regrouped families just
defined, including the empty family, and let
$A$ range over subsets of the coarse edges outside their
intervals. With the whole-contour scalar convention above,
the exact reconstruction is
\begin{equation}\label{eq:c8-fourth-coarse-family-identity}
    \frac{F^\Lambda(1)}{Z}
        =\frac{S^\circ_\Lambda}{Z}
            \sum_{\mathfrak D}
                q^\circ_\Lambda
                    \left(\bigcup_{\mathcal D\in\mathfrak D}\mathcal D\right)
                \sum_A
                    \mathcal C_V
                        ((g_{\mathcal D})_{\mathcal D\in\mathfrak D},
                         (\ell_a)_{a\in A})\;.
\end{equation}
For an empty union the reference factor is one. For both
families empty, the contraction is the coarse Wilson character.
This identity has one global factor $Z^{-1}$, the original
joint exclusion weight, and the original scalar factor.
For a finite cutoff with $N_c$ coarse traversals, the pin
input bounds $\sum_W b_W$ by $LN_c\epsilon$. Hence the
positive fine-family majorant is at most $e^{LN_c\epsilon}$.
Expanding the remaining coarse gaps costs at most
$\prod_a(1+\|\ell_a\|_\star)\leq3^{N_c}$.
These finite-volume majorants justify the regrouping at this
scope. They are not used as uniform perimeter bounds.

\paragraph{A quantitative pin bound for the merged source blocks.}
Put $b_W=e^{\theta^\circ(W)}\|r_W\|_\star$ and
$C_I=M+2$, where $LM$ is the unit used in
$\sigma(W)=n(W)+|\mathcal H(W)|/(LM)$. Since each coarse
traversal has $L$ fine traversals, the literal pin input gives,
for every coarse traversal $c$,
\begin{equation}\label{eq:c8-fourth-coarse-pin-input}
    \sum_{W:\,c\in I_W}b_We^{2\nu\sigma(W)}
        \leq L\epsilon,\qquad
    |I_W|\leq M\sigma(W)+2\leq C_I\sigma(W)\;.
\end{equation}
The first inequality counts at most the $L$ fine positions
in that coarse edge; every $W$ whose rounded interval contains
it meets at least one of those positions.
Thus, for every $W$,
\begin{equation}\label{eq:c8-fourth-overlap-budget}
    \sum_{W':\,I_{W'}\cap I_W\neq\varnothing}
        b_{W'}e^{2\nu\sigma(W')}
        \leq C_I L\epsilon\,\sigma(W)\;.
\end{equation}
Assume $C_I L\epsilon\leq\nu$. Apply the existing rooted
connected-tree bound with node weight $b_We^{\nu\sigma(W)}$
and node cost $\nu\sigma(W)$ to this actual overlap relation.
Its required row sum is exactly
\eqref{eq:c8-fourth-overlap-budget}. The bound for a fixed
root is then $b_We^{2\nu\sigma(W)}$. More explicitly, the
connected-set sum is bounded by the rooted-tree sum with the
factor $1/(j-1)!$ for the ordered other $j-1$ vertices.
Allowing repeated vertices only enlarges this positive bound;
the physical families still have their original distinct
group labels and compatibility restrictions.
The corresponding positive rooted-tree recursion starts at
$T_W^{(0)}=b_We^{\nu\sigma(W)}$ and updates by
$T_W^{(j+1)}=b_We^{\nu\sigma(W)}
\exp(\sum_{W':\,I_{W'}\cap I_W\neq\varnothing}T_{W'}^{(j)})$.
Substitution of \eqref{eq:c8-fourth-overlap-budget} proves
$T_W^{(j)}\leq b_We^{2\nu\sigma(W)}$ at every depth.
This is the specific reserve payment used here.

Every merged interval containing $c$ has at least one original
$W$ with $c\in I_W$. Rooting there, using
\eqref{eq:c8-fourth-merged-norm}, and applying the pin bound
gives
\begin{equation}\label{eq:c8-fourth-merged-pin}
    \sum_{\mathcal D:\,c\in I_{\mathcal D}}
        e^{\nu s(\mathcal D)+
            \sum_{W\in\mathcal D}\theta^\circ(W)}
        \|g_{\mathcal D}\|_\star
        \leq L\epsilon\;.
\end{equation}
For a whole-contour component its scalar modulus replaces
the norm. The sum is over overlap-connected original
compatible sets. Overcounting a set by its possible roots
is harmless for this upper bound.
This is a bound on a sum, not a fixed-support exponential
estimate. No independence of the retained physical fields
has been used.

The extra merging cost is therefore paid by the specified
first-step reserve once $C_I L\epsilon\leq\nu$; it introduces
no global perimeter multiplier. The separate zero-root
species have the row-curvature cost
\eqref{eq:c8-fourth-zero-root-bound}. Their local smallness,
and the required strengthened deviation estimate for the
merged blocks, have not been inferred from
\eqref{eq:c8-fourth-merged-pin} alone.

\paragraph{Physical use and what is still missing.}
Substitute \eqref{eq:c8-fourth-coarse-family-identity} in the
same first-step numerator containing $F^\Lambda(1)$, with
all original labelled weights, Jacobians, cutoffs and pending
integrations. Divide by the same physical $D_h$.
The reference scalar $S^\circ_\Lambda$ is kept in that
numerator; it is not a physical normaliser. The current
boundary factors may then multiply the entire reconstructed
closed source before any inherited integration that reads
them. In particular neither the dependence on $y$ nor the
joint factor $q^\circ_\Lambda$ is averaged separately.

This completes the specified real, first-step family
regrouping conditional on its exact V3 inputs. It does not
make V3's theorem an all-scale input. Full read-out supports
in coarse variables, the uncovered data outside the supplied
scope, complex backgrounds, the next-step integration of
these closed kernels, and the uniform physical estimate
relative to $D_h$ still require their own bounds. The
zero-root curvature kernels and the global factor
$S^\circ_\Lambda/Z$ must remain in that continuation.
This subsection treats one contour. For the joint insertion of
two disjoint loops, selected components reaching both contours remain in
the original joint expansion. They cannot be replaced by a
product of the two one-loop conditional expectations or
discarded when applying the pin estimate.

\subsection{A Selected Integration Domain Reaching Both Disjoint Loops}
\label{sec:c8-fourth-disjoint-root}

The next operand is genuinely joint: one selected integration
domain with probes reading both disjoint loops. This calculation
does not multiply separately integrated one-loop images. It
estimates that contribution to the original two-loop insertion
and its touching vacuum-denominator completion. Chains of
several selected domains, and the full connected covariance,
are not identified with this one contribution.

\paragraph{Sources and the common selected measure.}
Fix the actual first-step retained data and carrier of V3.
For $i=1,2$, let $f_i(B;H_i)$ be the normalised closed character
of loop $C_i$ in the reconstructed original fine field at the
young variable $m+B$, with a fixed auxiliary closed $H_i$ at
its base, in the fixed unitary representation $R_i$ of dimension
$d_i$. Thus
\begin{equation}\label{eq:c8-fourth-two-loop-source}
    f_i(B;H_i)=\frac{1}{d_i}\operatorname{Tr}\left[
        R_i(\operatorname{Hol}_{U_\xi(m+B)}C_i)R_i(H_i)
    \right]\;.
\end{equation}
The fine holonomy and $H_i$ are conjugated at the same base
under the induced gauge action, so each source is invariant.
The physical value has $H_i=1$. The contours $C_1$
and $C_2$ are disjoint. Both sources use the same reconstructed
field and the same retained data. Let $\mathsf S_i$ be their
field-independent sets of young free-bond inputs, including
every row input on which their reconstructed links depend.
This paragraph assumes $\mathsf S_1\cap\mathsf S_2=\varnothing$.
For the first-step local reconstruction this follows when
the two full row-readout neighbourhoods are disjoint. It does
not follow just from disjoint contours for an arbitrary
inherited source.

Use the actual split-selection measures
$\Lambda^{f,1}_{\Lambda;\ell}$ and the same final domains
$Z(\ell)$ of V3 Lemma \texttt{c5-5-2}. Each
$\ell$ retains its selected ordinary hosts, probe set
$\mathsf J$, bad cutoffs, enlarged seed, weakened integral
and original scalar normalisation. The inputs used here are
its dummy comparison and the closed-source finite-difference
bound of Lemma \texttt{c5-5-7z}. For every connected final
domain $Z$ and every nonempty probe set $J_i\subseteq\mathsf S_i$,
they read
\begin{equation}\label{eq:c8-fourth-joint-root-inputs}
    \begin{aligned}
        \sum_{\ell:\,Z(\ell)=Z}
            u_*^{|\mathsf J(\ell)|}
                \|\Lambda^{f,1}_{\Lambda;\ell}\|_{\rm TV}
            &\leq m(Z)
                =\varepsilon_*e^{-\lambda_*\ell_1(Z)},\\
        \sup_{H_i,B}|\Delta_{J_i}f_i(B;H_i)|
            &\leq u_i^{|J_i|},\qquad
        u_i=c_{{\rm src},i} g_0(1+2r)\leq u_*\;.
    \end{aligned}
\end{equation}
The supremum is on the real protected box containing the
supports and all finite-difference evaluation points.
$c_{{\rm src},i}$ is the fixed character and row-reconstruction constant
for the square $C_i$, independent of its length. No Gaussian
measure is substituted for the displayed selected measure.
These first-step inputs retain their stated hypotheses;
their extension through arbitrary inherited integrations
is not assumed.

For a selection with probes in $\mathsf S_1\cup\mathsf S_2$,
put $\mathsf J_i=\mathsf J\cap\mathsf S_i$. The disjoint
read-out sets imply the exact finite-difference identity
\begin{equation}\label{eq:c8-fourth-two-source-difference}
    \Delta_{\mathsf J}(f_1f_2)
        =(\Delta_{\mathsf J_1}f_1)
            (\Delta_{\mathsf J_2}f_2)\;.
\end{equation}
This is factorisation of a function under differences,
not factorisation of its integral. A selected probe outside
both read-out sets gives zero.
For selections with both $\mathsf J_i$ nonempty define
\begin{equation}\label{eq:c8-fourth-joint-root-integral}
    A_\ell^{12}(H_1,H_2)
        =\frac{1}{Z_1Z_2}
            \int
                (\Delta_{\mathsf J_1}f_1(B;H_1))
                (\Delta_{\mathsf J_2}f_2(B;H_2))\,
                    \dd\Lambda^{f,1}_{\Lambda;\ell}(B)\;.
\end{equation}
This is one product of two complete closed contractions
inside the same integration. The scalars $Z_1,Z_2$ are the
prescribed source factors, each occurring once.
At $H_1=H_2=1$, the C8 Boolean identity expresses
$f_i/Z_i$ as the sum of its actual geometric marked
coefficients. The empty geometric coefficient is the
coarse-loop character divided by $Z_i$, independent of
the young variables. Since each $\mathsf J_i$ is nonempty,
its finite difference kills that empty coefficient.
Thus \eqref{eq:c8-fourth-joint-root-integral} is the same
selected integral of the signed sum of C8 terms with
nonempty geometric marked sets on both loops. The geometric
marked sets and the free-bond probe sets $\mathsf J_i$
are different labels. Both geometric sums occur inside the
same integral, before any estimate.

Put $\rho_i=u_i/u_*$. Since both $|\mathsf J_i|\geq1$ and
$\rho_i\leq1$, the two source gains survive the joint
integration and the full split-selection sum. Write
$\mathscr L_{12}(Z)$ for the selections with final domain $Z$,
$\mathsf J\subseteq\mathsf S_1\cup\mathsf S_2$, and both
$\mathsf J_i$ nonempty. Then
\begin{equation}\label{eq:c8-fourth-joint-root-bound}
    \sum_{\ell\in\mathscr L_{12}(Z)}
        \sup_{H_1,H_2}|A_\ell^{12}(H_1,H_2)|
        \leq
            \frac{\rho_1\rho_2}{|Z_1Z_2|}
                \varepsilon_*e^{-\lambda_*\ell_1(Z)}\;.
\end{equation}
Indeed $u_1^{|\mathsf J_1|}u_2^{|\mathsf J_2|}
\leq\rho_1\rho_2u_*^{|\mathsf J|}$, after which the first
line of \eqref{eq:c8-fourth-joint-root-inputs} applies.
No independence between the fields near the disjoint loops
is used.

\paragraph{The touching denominator labels.}
For this single selected domain, the original conditional
vacuum ratio is, with $\operatorname{Hit}(Z)$ denoting its
touching vacuum clusters,
$q_\Lambda^v(Z)=\exp(-\sum_{\mathfrak c\in\operatorname{Hit}(Z)}
\Phi_\Lambda^{v,1}(\mathfrak c))$, as in V3
Corollary \texttt{c5-4-3}. Denote a term of its exponential
expansion by an ordered cluster list $\boldsymbol{\mathfrak c}$,
with its factor $(-1)^j/j!$. Its source coefficient is
$A_\ell^{12}$ times that term; call it
$F_{\ell,\boldsymbol{\mathfrak c}}^{12}$.
Repeated cluster labels and their internal connected Ursell
sums remain exactly as in the vacuum expansion.
The actual hitting bound is
$\sum_{\mathfrak c\in\operatorname{Hit}(Z)}
|\Phi_\Lambda^{v,1}(\mathfrak c)|\leq N_1(Z)/8$.
Taking absolute values only after the connected graph sums
therefore gives
\begin{equation}\label{eq:c8-fourth-joint-root-completed}
    \sum_{\ell\in\mathscr L_{12}(Z)}
        \sum_{\boldsymbol{\mathfrak c}}
        \sup_{H_1,H_2}
            |F_{\ell,\boldsymbol{\mathfrak c}}^{12}|
        \leq
        \frac{\rho_1\rho_2\varepsilon_*}{|Z_1Z_2|}
            e^{-(\lambda_*-C_{\rm g}/8)\ell_1(Z)}\;.
\end{equation}
The selection sum here has exactly the restrictions in
\eqref{eq:c8-fourth-joint-root-bound}. The inequality uses
$N_1(Z)\leq C_{\rm g}\ell_1(Z)$, with $C_{\rm g}=162$,
and sums the complete touching-cluster series by its
exponential majorant. Thus it is an absolute sum over these
denominator labels, not only a bound for their signed sum.
Their supports are not erased; no fixed-support bound on
the full completed source is inferred.

\paragraph{The separation payment.}
Let $\mathcal P_i$ be the set of $LM$-cubes containing the
initial vertices of bonds in $\mathsf S_i$, and let
$n_i^{\rm pin}=|\mathcal P_i|$.
If either set is empty, the contribution above vanishes.
Otherwise define
\begin{equation}\label{eq:c8-fourth-source-separation}
    \Delta_{12}=
        \frac{1}{LM}
            \min_{\Delta_i\in\mathcal P_i}
                \operatorname{dist}(\Delta_1,\Delta_2)\;,
\end{equation}
where distance is the periodic Euclidean distance between
the closed cubes. A domain carrying a probe from each loop
contains a cube from each $\mathcal P_i$. The connected tree
in the definition of $d_1(Z)$ meets both cubes, so
$\ell_1(Z)=d_1(Z)+5\geq\Delta_{12}$.
Put $\lambda_c=\lambda_*-C_{\rm g}/8$ and $s=\lambda_c/2$.
The native choice $\lambda_*\geq2C_{\rm g}\log128$ gives
$s>C_{\rm g}\log64$. V3's domain-counting bound is therefore
available with
$B(s)=e^{-s/C_{\rm g}}/(1-64e^{-s/C_{\rm g}})$.
Splitting the exponential between the span and the summation
and pinning a cube on the smaller source set gives the following
bound. Here $\mathscr Z_{12}$ is the set of final domains
carrying a probe of each source, and the selection sum uses
$\mathscr L_{12}(Z)$.
\begin{equation}\label{eq:c8-fourth-disjoint-root-tail}
    \sum_{Z\in\mathscr Z_{12}}
        \sum_{\ell,\boldsymbol{\mathfrak c}}
            \sup_{H_1,H_2}
                |F_{\ell,\boldsymbol{\mathfrak c}}^{12}|
        \leq
        \frac{\rho_1\rho_2\varepsilon_*}{|Z_1Z_2|}
            \min(n_1^{\rm pin},n_2^{\rm pin})B(s)e^{-s\Delta_{12}}\;.
\end{equation}
For the last summation, each admissible $Z$ contains at least
one cube of either chosen $\mathcal P_i$. Sum
$e^{-s\ell_1(Z)}$ over domains containing that cube and use
the animal bound; overcounting domains is harmless.
The remaining factor $e^{-s\ell_1(Z)}$ pays the separation.
This is decay in the distance between the full young-variable
read-out sets, not a claim that two disjoint contours alone
give independent integrations.

\paragraph{A uniform first-step bound and its positive inherited lift.}
The scalar consequences in this paragraph and the following
fourth-response paragraph are diagnostic only. They do not
propagate a local labelled marked-source norm and are not
used as an induction step of the active marked-polymer argument.
For a bound uniform over the retained first-step labels, let
$R_{{\rm src},i}$ bound the distance from a free input in a
source link's stencil to that link, and let $m_{{\rm src},i}$
bound the number of free inputs in such a stencil. These are
the fixed row-stencil constants of the actual first-step
reconstruction. Let $n_{C_i}$ be the number of fine traversals.
Suppose the distance between the contours exceeds
$R_{{\rm src},1}+R_{{\rm src},2}$, so the two read-out sets
are disjoint on every retained label in the stated scope.
This condition concerns the first-step readouts; it holds
eventually along a refining lattice for fixed separated
physical loops, not at every coarse scale.
Set
\begin{equation}\label{eq:c8-fourth-uniform-disjoint-constant}
    \begin{aligned}
        \Delta_*&=
            \frac{(\operatorname{dist}(C_1,C_2)
                -R_{{\rm src},1}-R_{{\rm src},2}-4LM)_+}{LM},\\
        P_*&=\min(m_{{\rm src},1}n_{C_1},
                    m_{{\rm src},2}n_{C_2}),\\
        C_{12}^{\rm one}
            &=\frac{\rho_1\rho_2\varepsilon_*}{|Z_1Z_2|}
                P_*B(s)e^{-s\Delta_*}\;.
    \end{aligned}
\end{equation}
Each $LM$-cube in four dimensions has diameter $2LM$.
Consequently $\Delta_{12}\geq\Delta_*$, while the stencil
count gives $n_i^{\rm pin}\leq m_{{\rm src},i}n_{C_i}$.
These facts give the stated constant without an unquantified
loop-length factor.

Let $B_{12}^{\rm one}(y)$ be the sum in
\eqref{eq:c8-fourth-disjoint-root-tail}, with its original
signs, at $H_1=H_2=1$, multiplied by the original labelled
data indicator $\mathbf 1_{\mathscr E}$ on which its inputs
hold. Set this selected contribution equal to zero off that
data set; the complementary contribution to the original
source is retained separately. The source remains a sum
of invariant closed contractions with the complete original
gauge data. The estimate proves
\begin{equation}\label{eq:c8-fourth-initial-connecting-source}
    |B_{12}^{\rm one}(y)|\leq C_{12}^{\rm one}\;.
\end{equation}
This is the signed sum of the specified single-root
contributions and their denominator labels, not the full
joint Wilson source.

Now propagate this same contribution by the actual real
chronological lift WC-C8.3a--WC-C8.3b. At a step with
nonnegative source-independent vacuum weight $w_\lambda$,
write its original normaliser as $I(y)=\sum_\lambda
\int w_\lambda(z,y)\,\dd z$. Before division, the bound is
\begin{equation}\label{eq:c8-fourth-connecting-lift}
    \left|\sum_\lambda\int
        w_\lambda(z,y)B_{\rm previous}(z,y)\,\dd z\right|
        \leq C_{12}^{\rm one}I(y)\;.
\end{equation}
Thus the same constant bounds the next quotient on its
nonzero-mass outputs. In an R replacement the common factor
$\sigma(v,y)/J(y)$ multiplies both
$\int r(w,y)B_{\rm previous}(w,y)\,\dd w$ and
$\int r(w,y)\,\dd w$; nonnegativity of the original real
$r,\sigma$ gives the same conclusion. Null outputs carry
zero physical mass. Postponed variables are kept in $y$.
This is the positive-lift argument for the actual vacuum
weights, including all old integrations and normalisers,
not a replacement law. It keeps the same prescribed
$Z_1,Z_2$ already present in
$C_{12}^{\rm one}$.

Accordingly the inherited, unlabelled contribution $B_h^{\rm one}$
has the same bound through any sequence of these positive
real lifts. It is kept as a whole contribution with its full
read-out support. No bound on the absolute values of new
labels created by a later localisation is inferred.
If $\mathcal I_h$ denotes the full physical vacuum integration,
with $\mathcal I_h[1]=D_h>0$, then
\begin{equation}\label{eq:c8-fourth-connecting-physical}
    \frac{|\mathcal I_h[B_h^{\rm one}]|}{D_h}
        \leq C_{12}^{\rm one}\;.
\end{equation}
The numerator is bounded by $C_{12}^{\rm one}\mathcal I_h[1]$
under the original positive physical law. This uses its
actual normalisation; it does not replace $D_h$ by a
conditional component normaliser.

\paragraph{Fourth boundary response of this inherited contribution.}
Use the same bounded boundary tests as in
\eqref{eq:c8-fourth-defined-test}, now bounded on the whole
physical space of $\mathcal I_h$. Define the diagnostic
response by
\begin{equation}\label{eq:c8-fourth-connecting-response}
    F_{{\rm one},h}^{ij}(s)
        =\frac{\mathcal I_h[
            B_h^{\rm one}e^{s_1p_i+s_2q_i+s_3p_j+s_4q_j}]}
            {\mathcal I_h[
                e^{s_1p_i+s_2q_i+s_3p_j+s_4q_j}]}\;.
\end{equation}
The entire inherited contribution and the physical denominator
are retained. The source is fixed during this current
boundary deformation, exactly as in the earlier diagnostic.
Apply the cumulant formula
\eqref{eq:c8-fourth-cumulant} with
the bound \eqref{eq:c8-fourth-initial-connecting-source}
propagated by \eqref{eq:c8-fourth-connecting-lift}.
Every source-containing block costs at most
$C_{12}^{\rm one}$ times its test suprema. The sum of the
five-label partition weights is $150$, so
\begin{equation}\label{eq:c8-fourth-connecting-fourth-bound}
    \left|\sum_{i,j}
        \partial_{[4]}F_{{\rm one},h}^{ij}(0)\right|
        \leq150\,C_{12}^{\rm one}
            \left(\sum_i
                \|p_i\|_\infty\|q_i\|_\infty\right)^2\;.
\end{equation}
Its complex boundary-source neighbourhood has the same
physical denominator lower bound as
\eqref{eq:c8-fourth-vacuum-lower}. The bound is for this
specified inherited contribution. It does not bound the
whole Wilson covariance, replace the inherited physical
law, or establish a noncommuting field-rotation identity.

\paragraph{Physical meaning and limits of this estimate.}
The coefficient is a term in the joint source, not yet the
full connected covariance. The latter additionally subtracts
the product of the separate conditional images, and includes
connections through several selected domains. Those terms
have not been bounded here.
The estimate may be inserted in the original physical
numerator at its stated data set, with all outer weights and
the same $D_h$. The prescribed $|Z_1Z_2|^{-1}$ remains in
the bound. An estimate uniform in the ultraviolet cutoff for
the full source still has to pay it and all other contributions;
no ultraviolet conclusion
is obtained just by writing the separation exponential.

In particular a later inherited component $Y_i$ can read both
loops even when their contours are disjoint. Its full support
must then remain in the source. LFR2 (1.103) does not supply
tree decay inside the fixed $Y_i$ through its normalised
operation. The positive lift above preserves the already
obtained global bound for the selected initial contribution;
it does not give $e^{-\lambda_*\ell_1(Y_i)}$ for a new
inherited component or an absolute bound on all its newly
created localisation labels. Such an estimate still requires
the inherited source analysis and the physical outer weight.
Empty read-out sets give zero, and $\Delta_{12}=0$ gives no
separation gain. With one probe on each loop, the product of
their linear source increments is even under a common sign
reversal, so parity alone does not force this joint term
to vanish.

\subsection{Keeping the Geometric Marks in the Native Source Estimate}
\label{sec:c8-fourth-fixed-geometric}

This estimate keeps the geometric sets $J_1,J_2$ fixed and
takes absolute values separately in the native selection and
denominator labels. It addresses the marked-source factor
left unresolved by a bound on a global unlabelled sum.
The setting is the actual native first-step integral on real
retained data. A transfer of its hypotheses to later inherited
kernels is not assumed.

\paragraph{The raw input kernels and their norm.}
For each disjoint loop, let its nonempty $J_i$ lie in a specified
straight or corner interval $I_i$ to which WC-C8.7 applies.
Include all the interval, tree and return paths in its support.
Apply its local closing substitution to the raw WC-C8.3
insertion, with its internal physical holonomies restored,
and denote this raw input by $k_i(B;H_i)\in E_{\chi_i}$.
This section's $k_i$ is not the completed WC-C8.7 kernel
constructed from the inherited, integrated $f_{J,\sigma,X}$.
The external
$H_i$ remains a complete closed holonomy. All internal
coarse transports are frozen during the young-variable
integration. Define
\begin{equation}\label{eq:c8-fourth-elementary-mark}
    \kappa_{i,a}(B;H)=
        \chi_i(Q_{i,a}(B)H)-\chi_i(H)\;.
\end{equation}
These are complete closed functions, with the induced
simultaneous conjugation of $Q_{i,a}$ and $H$. In the
already defined algebra $E_{\chi_i}$, this raw input
kernel is
\begin{equation}\label{eq:c8-fourth-marked-star-word}
    k_i(B)=\frac{1}{Z_i}\,
        \kappa_{i,a_1}(B)\star u_{i,1}\star\cdots
        \star\kappa_{i,a_{m_i}}(B)\star u_{i,m_i}\;.
\end{equation}
Here the $u_{i,r}$ are the closed kernels of the actual
internal physical holonomies and the final $L_{\rm loc}$.
They have norm one at real data and are independent of $B$.
This identity is just the original ordered complete trace;
no product of its individual characters is substituted.
For the pair use the fixed tensor-product closed-kernel
space $E_{\chi_1}\otimes E_{\chi_2}$ and its left-multiplication
operator norm, denoted $\|\cdot\|_{\star,12}$.
For elementary tensors this norm is multiplicative.
Its auxiliary Haar construction is an algebraic norm,
not an independence assumption for the physical fields.

\paragraph{Derivative and undifferentiated-factor inputs.}
Let $\mathsf F_{i,a}$ be the young free-bond readout of the
fine arc belonging to $a$, including all row inputs.
At this first step each such arc has $L$ fine traversals.
The increment variables lie in the real box
$|B_b|_\sigma\leq5r/4$ on the native young carrier $D(\Lambda)$;
the retained real shift satisfies $|m_b|_\sigma\leq r/4$.
Zeroing or scaling any selected increment therefore keeps
$m+B$ in V3's enlarged box $|m_b+B_b|_\sigma\leq2r$.
The same rows, cutoffs and retained fields are used.
The analytic row estimates imply
the following bound for distinct free-bond labels in
$\mathsf F_{i,a}$ and unit Lie-algebra directions:
\begin{equation}\label{eq:c8-fourth-mark-derivative}
    \sup_B
        \|\partial_{e_1,b_1}\cdots\partial_{e_q,b_q}
            \kappa_{i,a}(B)\|_{\star,i}
        \leq(C_{\partial,i}g_0)^q\;.
\end{equation}
Here $1\leq q\leq|\mathsf F_{i,a}|$. A derivative in a label
outside that readout is zero. The constant is independent
of $J_i$, the loop length and the torus. To obtain it, apply
the product rule inside the closed character of this
$L$-bond arc and the analytic row derivative bounds on
the enlarged phase chart. Only finitely many derivative
orders occur on a row. The finite-dimensional inequality
$\|f\|_{\star,i}\leq K_{\chi_i}\sup_{H\in G}|f(H)|$
absorbs its constant into $C_{\partial,i}^q$ for $q\geq1$.
The gauge data and the derivative directions transform
together; only norms of complete closed kernels enter.
No assertion with a constant independent of depth is made
for an original fine arc of length $L^k$ at a later scale.

Partition the geometric marks into good and remaining marks.
For each good mark suppose
$\sup_B\|\kappa_{i,a}(B)\|_{\star,i}\leq\varepsilon$,
where $0<\varepsilon\leq1$ is common to the two loops.
Every remaining mark retains the unconditional real bound
$\sup_B\|\kappa_{i,a}(B)\|_{\star,i}\leq2$.
One sufficient check for a good mark is explicit:
if its base row satisfies the hypotheses of
\eqref{eq:c8-fourth-coarse-curvature}, then
\begin{equation}\label{eq:c8-fourth-good-mark-check}
    \sup_B\|\kappa_{i,a}(B)\|_{\star,i}
        \leq
        2C_{R_i}D^2L^2\delta_a(U_\xi(m))
        +\frac{5r}{4}|\mathsf F_{i,a}|C_{\partial,i}g_0\;.
\end{equation}
The first term bounds the actual base relative holonomy;
the second follows by integrating
\eqref{eq:c8-fourth-mark-derivative} along the straight
line from the zero young increment to $B$.
The mean shift and retained exterior data stay in the base
field. A mark is not declared good without this or another
bound on its actual kernel. Thus no large-field factor is
hidden in $\varepsilon$.

\paragraph{Paying both the marks and the probes.}
Let $m_{\rm g}$ and $m_{\rm b}$ be the numbers of good and
remaining geometric marks in the pair. Let $N_{\rm inc}$
bound the number of marked-arc readouts containing any one
young free label, counting both loops. The first-step row
locality and square traversal multiplicity give such a
constant independent of their lengths. Put
\begin{equation}\label{eq:c8-fourth-grade-parameters}
    \eta=\sqrt{\varepsilon},\qquad
    v=N_{\rm inc}\max_i C_{\partial,i}g_0(1+2r),\qquad
    \widetilde u=\frac{v}{\eta},\qquad
    \widetilde\rho=\frac{\widetilde u}{u_*},\qquad
    \mathfrak a_J=\frac{\eta^{m_{\rm g}}2^{m_{\rm b}}}{|Z_1Z_2|}\;.
\end{equation}
Require $\widetilde u\leq u_*$. For instance,
$v\leq\varepsilon\leq1$ and $v\leq u_*^2$ imply this
condition. Let $\mathsf J$ be a set of distinct young probes,
and put $q=|\mathsf J|$. The finite-difference formula gives
\begin{equation}\label{eq:c8-fourth-fixed-mark-difference}
    \sup_B\|
        \Delta_{\mathsf J}(k_1\otimes k_2)(B)
    \|_{\star,12}
        \leq\mathfrak a_J\,\widetilde u^q\;.
\end{equation}
Here $\Delta_{\mathsf J}$ acts on the actual joint kernel
with the same young field in both factors. Disjoint
physical contours do not have to have disjoint readouts
for this bound.

For the proof, apply the product rule inside the ordered
closed-kernel products in
\eqref{eq:c8-fourth-marked-star-word}.
Assign each derivative label to a marked factor reading it;
there are at most $N_{\rm inc}^q$ assignments.
An undifferentiated good factor costs
$\varepsilon=\eta^2\leq\eta$. If it receives $j\geq1$
derivatives, its increment bound from
\eqref{eq:c8-fourth-mark-derivative} is no larger than
$\eta(C_{\partial,i}g_0(1+2r)/\eta)^j$,
because $\eta^{1-j}\geq1$.
An undifferentiated remaining factor costs two; its
differentiated bound is also dominated by two times the
same probe factor. Multiplying the factors, summing the
assignments and integrating the unit-volume
finite-difference interpolation cube proves the claim.
There is no factorial in the total number of probes.
In the dangerous case of one good mark and one derivative,
the allocated factors give $\eta(v/\eta)=v$, not the false
extra gain $\varepsilon v$.

\paragraph{The labelled native caller.}
Insert this same joint kernel into one actual selected
measure and retain its geometric labels:
\begin{equation}\label{eq:c8-fourth-fixed-mark-selected}
    K_{\ell;J_1,J_2}
        =\int
            \Delta_{\mathsf J(\ell)}(k_1\otimes k_2)(B)\,
                \dd\Lambda^{f,1}_{\Lambda;\ell}(B)\;.
\end{equation}
Both closed kernels remain inside the common integration;
no product of separately integrated kernels is used.
Fix a positive integer $q_0$ and consider exactly the native
selections with at least $q_0$ actual probes. For every final
domain $Z$, the dummy comparison gives
\begin{equation}\label{eq:c8-fourth-fixed-mark-root}
    \sum_{\substack{\ell:\,Z(\ell)=Z\\
        |\mathsf J(\ell)|\geq q_0}}
        \|K_{\ell;J_1,J_2}\|_{\star,12}
        \leq
        \mathfrak a_J\,\widetilde\rho^{q_0}
            \varepsilon_*e^{-\lambda_*\ell_1(Z)}\;.
\end{equation}
The proof is
$\widetilde u^{|\mathsf J|}
\leq\widetilde\rho^{q_0}u_*^{|\mathsf J|}$
followed by the original selection sum.
For disjoint readouts and a probe from each loop one may
take $q_0=2$. Mere disjointness of the contours does not
justify that choice.
The sector with no probes is retained separately; its
base-value kernel has the direct bound
$\varepsilon^{m_{\rm g}}2^{m_{\rm b}}/|Z_1Z_2|$.
It is not removed by an estimate with $q_0\geq1$.

Completing this source with its original touching
vacuum-denominator cluster labels, as above, replaces the
last exponential by
$e^{-(\lambda_*-C_{\rm g}/8)\ell_1(Z)}$ in the corresponding
absolute sum of these operator norms. The factor
$\eta^{m_{\rm g}}2^{m_{\rm b}}$ and the original
$|Z_1Z_2|^{-1}$ remain. In particular the geometric marks
have not been summed with signs and then forgotten.
The complete supports retain both germ intervals, all
coarse and return transports, every reconstructed-field
readout and all denominator clusters.

This is a bound in the existing closed-kernel norm for
specified labelled native contributions. It is not yet an
estimate with a tree weight on their full enlarged supports.
Paying the remaining-mark factor $2^{m_{\rm b}}$, summing
the actual marked families and inherited histories, and
proving the same hypotheses on the required complex
domains are still necessary. Those estimates cannot be
replaced by the global positive-lift diagnostic above.
The full physical denominator and the prescribed scalar
extractions remain unchanged at the downstream caller.

\paragraph{Compact checkpoint: C8-FO-1.}
The target is the actual C8 insertion for disjoint loops, under
the full inherited law and physical denominator $D_h$.
WC.1 and the NLSM-style noncommuting rotation inequality remain
open. The calculation stays in this separate record; the
original first-step file and V3 are untouched.

The conditional fourth-response bound and its physical
denominator neighbourhood were checked in
\path{review-fo3/report.md}. The occurrence-reserve estimate
\eqref{eq:c8-fourth-labelled-bound} was checked in
\path{review-fo4/report.md}; that review correctly says it
does not itself estimate the inherited source.

The local C8 regrouping and real coarse-word curvature bound
were checked as conditional first-step implications in
\path{review-fo5/report.md}. The simultaneous first-step
reconstruction \eqref{eq:c8-fourth-coarse-family-identity}
and merged pin \eqref{eq:c8-fourth-merged-pin} were checked
conditionally in \path{review-fo6/report.md}, with its
\path{review-fo6/disclosure-supplement.md}. The V3 foundations
were explicit inputs to these reviews, not independently
certified. These one-loop preparatory results do not close
the joint inherited caller.

The current new joint calculation is
\eqref{eq:c8-fourth-disjoint-root-tail}, for one selected domain
carrying a probe from each disjoint loop, with all touching
vacuum-denominator labels. Its pointwise bound is carried
through the actual positive real lifts in
\eqref{eq:c8-fourth-connecting-lift}, retaining the physical
denominator in \eqref{eq:c8-fourth-connecting-physical} and
giving the fourth response
\eqref{eq:c8-fourth-connecting-fourth-bound}.
Its conditional scalar conclusions were checked in
\path{review-fo7/report.md}. Objection FO7-SCOPE remains
open: \eqref{eq:c8-fourth-connecting-lift} preserves a global
unlabelled bound, not the local labelled marked norm needed
by WC-C8.6. That part is diagnostic only and is excluded
from the active induction. It does not close the full
connected covariance or labelled source estimate.
The complementary data, connections through several selected
domains, new inherited components meeting both loops,
complex-background control and the prescribed inverse
$Z_1Z_2$ still require estimates.

The latest IT-1 and D2 records have been checked for relevant
changes. D2.2 now contains the next-step substitution and
gauge-covariance discussion; its coupled source expansion
and integration are not treated here as established inputs.
The current repair keeps the C8 geometric sets fixed.
Equations \eqref{eq:c8-fourth-fixed-mark-difference} and
\eqref{eq:c8-fourth-fixed-mark-root} estimate the native
closed-kernel norm with an explicit mark/probe trade-off,
retaining the factor for remaining marks and all denominator
labels. The focused review is in \path{review-fo8/report.md}.
It validates the calculation conditionally for the raw native
input, not its application to the completed inherited kernel.
Objection FO8-OPERAND remains open at
\eqref{eq:c8-fourth-marked-star-word}: the reviewed revision
\path{review-fo8/argument-snapshot.tex}, SHA-256
\path{1930678b59cb1d0847b80cb7bb1c2a02a1dcdc9f3b7ec183b95dd89641e86c3b},
identified the raw star word with the completed WC-C8.7
kernel without proving that identification. The text now
withdraws that identification. Closure requires the actual
joint inherited source, its integrations and denominator
labels to satisfy the stated marked norm estimate. A product
of separately completed one-loop kernels does not supply
that joint operand.
The full-support weights, remaining-mark cost, inherited
field maps and complex domains are not certified by it.
The next calculation must address those actual marked
contributions, not a replacement global unlabelled carrier.
Frozen closed-kernel algebra,
the source definition and the exact physical law are not
to be rederived.

The newer D2.2(5)(d) also requires a spectator-dependence
check before its proposed self-energy can be used as a
single-polymer activity: its cluster sum is restricted to
clusters hitting exactly one member of the current family,
so that family remains an argument. The familywise ratio
identity is not challenged by this observation. The relevant
source is \path{WILSON_IT1_D2_STEP_MAP_2026_10_22.md},
lines 486--487, at SHA-256
\path{f1972456c36262da815471a12e90397c507658a33d77c93ec9dbb5d26bb4b869}.
The proposed factorisation repair was not carried out before
the stopping condition below. The newer draft and V3 are
unchanged by this record.

The user authorised sustained work until an observed usage
reset. Baseline on 22 October: 34 percent, primary
10080-minute window, reset timestamp 1793098031.
The preceding reading was 66 percent with that timestamp.
On 22 October, immediately before 17:26:21 UTC, the account
reported 2 percent and reset timestamp 1793293920 in the
same 10080-minute window. This is the observed reset.
Mathematical work is stopped and the continuation automation
is paused. Only the review qualifications and this checkpoint
were recorded after observing the reset. A new human
instruction is required to resume.

\begin{thebibliography}{99}
    \bibitem{OpenAIRotation}
    OpenAI, \emph{Exponential decay in two-dimensional classical
    $O(n)$ models}, 23 September 2026, Section~2,
    \texttt{sections/rotation.tex}, Git commit
    \texttt{1581d4cbdd69a0faf34b2089614fe2a814316a6a},
    \url{https://github.com/openai/math}.
\end{thebibliography}
\end{document}
